\documentclass[11pt]{article}
\usepackage[T1]{fontenc}
\usepackage[utf8]{inputenc}
\usepackage{lmodern,microtype}
\usepackage[margin=1in]{geometry}
\usepackage{amsmath,amssymb,amsthm,mathtools}
\usepackage{booktabs,longtable,array}
\usepackage[dvipsnames]{xcolor}
\usepackage{hyperref}
\hypersetup{colorlinks=true,linkcolor=MidnightBlue,citecolor=MidnightBlue,
urlcolor=MidnightBlue,pdftitle={Evolution equations and large-height heat asymptotics},
pdfauthor={Drafted for Edward Baker}}
\setlength{\emergencystretch}{2em}
\setlength{\parskip}{.35em}
\newcommand{\R}{\mathbb R}
\newcommand{\C}{\mathbb C}
\newcommand{\dd}{\,\mathrm d}
\newcommand{\epsg}{\epsilon_g}
\newcommand{\Mc}{\mathcal M}
\newcommand{\Sc}{\mathcal S}
\newcommand{\Gc}{\mathcal G}
\newcommand{\Acal}{\mathcal A}
\newcommand{\Log}{\operatorname{Log}}
\newcommand{\sech}{\operatorname{sech}}
\newcommand{\arsinh}{\operatorname{arsinh}}
\newtheorem{proposition}{Proposition}[section]
\newtheorem{remark}[proposition]{Remark}
\numberwithin{equation}{section}
\title{Evolution equations and\\large-height heat asymptotics\\
\large Controlled ans\"atze, arithmetic modes, and physical density}
\author{Drafted for Edward Baker}
\date{25 October 2026\\[.3em]\small Working research draft}
\begin{document}
\maketitle
\begin{abstract}
We organize large-height asymptotics for the canonical theta heat field
around its evolution equation, the residual of an approximate solution,
and the equation satisfied by the approximation error. The objective is
to understand and bound the field for every real heat time: first at
fixed time, then uniformly on compact time intervals, and separately in
joint time--height limits. Amplitude and phase equations generate
corrections, while canonical initial data, arithmetic remainders and
stability estimates determine whether those corrections approximate
the actual field. Negative time is forward diffusion and makes the
arithmetic modes' damping and phase drift explicit.

The existing bounded results are retained: a corrected positive-time
carrier on \([32,33]\), a buffered signed crossover through zero,
two distinct accuracy scales on \(1\le -t\le2\), and a chartwise
join on \([-2,0]\). A first evaluated evolution correction, retaining
one further Riemann--Siegel term and a bounded correction integral,
improves the fixed-negative-time field, spatial jets through three and
first time error to \(O(|G|X^{-5/4})\). Its actual residual is
\(O(|G|X^{-7/4})\); the remaining outer-relative allowance still grows
at fixed negative time. The exact-gamma model now extends uniformly
to every declared fixed compact negative interval bounded away from
zero, with explicit endpoint factors and sufficient onset gates.
Those factors deteriorate toward zero or infinite time. We display
the forcing that an ansatz residual
induces in the physical potential, density and current equations.
Improving and extending the field and derivative estimates is the
primary research task. The regulator transition near \(t=64\) remains
an important subsequent test of amplitude accuracy and observable
transfer. Exact identities, inherited estimates and proposed proof
methods are distinguished; no all-time asymptotic theorem or new
physical-transfer estimate is asserted.
\end{abstract}
\noindent\textbf{LLM acknowledgement.}
Drafted for Edward Baker with substantial LLM assistance. The opening
was prepared with GPT-6 Astra (\texttt{gpt-6-astra}), recorded effort
\texttt{max}, in parent chat0022. Notes03--06 used GPT-6.1 Sol
(\texttt{gpt-6.1-sol}), effort \texttt{ultra}, in ansatz chat0001;
Note07 summarizes that round. The evolution rewrite, Notes08--09,
used GPT-6 Astra, effort \texttt{max}, in chat0002. The correction and
compact-family results (Notes10/12, Reviews07/09), and their integrations
AN04/AN06 (Notes11/13), used GPT-6.1 Sol, effort \texttt{ultra}, in
chat0003. Their status remains proved in-note from named classical
inputs, with same-model internal analytical assessment only. Author
verification and independent specialist review remain necessary.
No numerical or machine-symbolic research was executed for this
integration; compilation and rendered-page checks are production evidence.
\newpage
\begingroup
\small
\setlength{\parskip}{0pt}
\tableofcontents
\endgroup
\newpage

\section{Objective, evidence, and the role of evolution}
\label{sec:objective}
The primary objective is to understand and bound the large-height
asymptotics of \(H(t,z)\) for all real times. The field formula should
expose amplitude, unwrapped phase and arithmetic structure. Its error
estimate may be elaborate, and different formulas may be appropriate
in different regions. The evolution equation is both a construction
method for these formulas and a way to track their errors. A useful
step must improve a stated approximation bound or extend a proved
domain; cancelling a formal residual alone is not such a step.

The program has three distinct time requirements:
\begin{enumerate}
\item for every fixed real \(t\), an asymptotic as \(X\to\infty\),
with the dependence of its constants and onset on \(t\) specified;
\item estimates uniform on prescribed compact time intervals,
including intervals crossing zero;
\item joint limits \(t=t(X)\), with explicit scaling regions,
possibly unbounded positive or negative time, and controlled overlaps.
\end{enumerate}
None follows by substituting a moving time into a fixed-time remainder.
The first evolution-guided improvement at intermediate negative times
is now proved at a bounded natural-scale scope in
Proposition~\ref{prop:first-negative-correction}. The exact-gamma
family is now controlled on every declared fixed compact negative
interval bounded away from zero in
Proposition~\ref{prop:compact-negative-family}. Through-zero assembly,
positive-time gaps and joint limits remain part of the wider objective.

The previous research round \cite{ansatz03,ansatz04,ansatz05,ansatz06}
established useful starting comparisons. It corrected the outer
positive-time carrier on \([32,33]\), retained the arithmetic cutoff
correction through zero, proved full buffered spatial and first time
errors for that crossover, and constructed a negative-time
representation that can be joined through zero on \([-2,0]\).
These are \emph{proved in-note from named classical inputs, with
same-model internal assessment only}. Their constants and onsets
are not numerically instantiated. Section~\ref{sec:map} records the
remaining coverage rather than treating those bounded domains as an
all-time result.

Evolution is already used in three substantive ways: the finite
negative-time representation solves the heat equation exactly on a
frozen chart; residual cancellation supplies a direct time-error
estimate; and a time join retains the derivative of its cutoff.
The first such evaluated arithmetic correction is now supplied by
Note10 \cite{ansatz10}; it pays a separate canonical initial-data
improvement and retains the full evolving arithmetic population.
Note12 \cite{ansatz12} makes the fixed compact-negative-time endpoint
and onset dependence explicit. Further corrections, through-zero
assembly and outer-scale resummation remain open. Section~\ref{sec:evolution} states the exact
error equation and the stability obligations that make this program
quantitative. Section~\ref{sec:observable-evolution} then records how
the residual enters the physical observables.

The parent manuscript \cite{parent} supplies a finite analytical
all-real-time/all-height construction on its declared strip, through
Notes54--55. Notes58--61, summarized in \cite{summary62}, develop
positive source compression, finite evaluation and filtering errors,
exact phase-index selection using onset freedom, and polynomial
additional-accuracy-bit cost at a fixed descriptor and query box.
Useful uniform height cost and the distinct full-chain audit remain
open. Those results provide source and transfer machinery, but do not
already give the short, quantitatively controlled asymptotic
description sought here. Their status is unchanged by this rewrite.

The transition near \(t=64\) belongs to the fixed physical regulator.
It remains important because it tests very fine amplitude accuracy
and the transfer from field to density. It follows the primary
field-asymptotic work in the present priority order. The
Gaussian-polynomial core and optional special-function simplifications
remain supporting models in the wider catalogue \cite{ansatz01}.

\section{Normalization, observables, and exact time anchors}
\label{sec:normalization}
Throughout,
\begin{align}
\Phi(u)&=\sum_{n\ge1}
 (2\pi^2n^4e^{9u}-3\pi n^2e^{5u})e^{-\pi n^2e^{4u}},
 \qquad u\ge0, \label{eq:source}\\
H(t,z)&=\int_0^\infty e^{tu^2}\Phi(u)\cos(zu)\dd u,
 \qquad H_t=-H_{zz}, \label{eq:heat}\\
H(0,z)&=\frac18\xi(s)=\Mc(s)\zeta(s),
 \qquad s=\frac{1+iz}{2}, \label{eq:anchor-zero}\\
\Mc(w)&=\frac{w(w-1)}{16}\pi^{-w/2}\Gamma(w/2).
 \label{eq:gamma}
\end{align}
The source has an even extension. The time variable \(t\) is heat time;
the usual zeta height is \(x/2\), not \(x\). Write \(z=x+iy\) and \(X=|x|\).
The displayed real-height formulas are analytically continued on each
separate buffered tail; \(|x|\) itself is not a holomorphic coordinate.

The fixed normalized spatial kernel is
\begin{equation}
k(v)=\frac{315}{256}(1-v^2)^4\mathbf1_{\{|v|\le1\}},\qquad
k_\ell(x)=\ell^{-1}k(x/\ell),\qquad K(x,y)=k_\ell(x)k_h(y),
\label{eq:kernel}
\end{equation}
with fixed positive widths \(\ell,h\).

For a holomorphic model \(F\), keep the same physical regulator
\(q=\epsg^2\) and fixed spatial kernel as for \(H\):
\begin{align}
S_F&=\log(|F|^2+q),&
\rho_F&=\frac{q|F_z|^2}{\pi(|F|^2+q)^2},&
R_F&=K*\rho_F, \label{eq:density}\\
U_F&=K*S_F,&
j_F&=\frac{qF_{zz}\overline{F_z}}{\pi(|F|^2+q)^2},&
J_F&=K*j_F. \label{eq:current}
\end{align}
The actual signed covariance is
\begin{equation}
Q_F=K*((S_F)_x^2-(S_F)_y^2)
       -((U_F)_x^2-(U_F)_y^2).
\label{eq:covariance}
\end{equation}
The variable-regulator terms in the parent evolution equations remain.
A convenient exact identity, away from zeros, is
\begin{equation}
\rho_F=\frac1{4\pi}\left|\frac{F_z}{F}\right|^2
 \sech^2\!\left(\log\frac{|F|}{\sqrt q}\right).
\label{eq:sech}
\end{equation}
At zeros use \eqref{eq:density}. Rescaling \(F\) without rescaling \(q\)
changes the density, so an accurate phase alone is not sufficient.

\subsection{Convolutions on the two sides of zero time}
For \(t>0\), the exact Gaussian identity \cite[Eq.~(35)]{polymath} is
\begin{equation}
H(t,z)=\frac1{\sqrt{\pi t}}\int_\R e^{-v^2/t}
                      \frac{\xi(s+v)}8\dd v.       \label{eq:positive-anchor}
\end{equation}
For \(T>0\), direct Gaussian integration in \eqref{eq:heat} gives
\begin{equation}
H(-T,z)=\frac1{\sqrt{4\pi T}}\int_\R e^{-v^2/(4T)}
                                      H(0,z-v)\dd v. \label{eq:negative-anchor}
\end{equation}
The even source and its rapid decay justify interchange on fixed
complex strips. These are exact starting identities, not short
large-height asymptotics. The negative-time formula must not be obtained
by substituting \(t<0\) into an uncontrolled positive-time remainder.

\section{Evolution of an ansatz and its error}
\label{sec:evolution}
This section integrates the exact identities in Ansatz Note08
\cite{ansatz08}. It provides a method and a checklist of estimates,
not a new remainder theorem.

\subsection{Separate the solution error from the residual}
For an ansatz \(F\), differentiable in time and holomorphic in space,
write
\begin{equation}
 \mathfrak r_F=F_t+F_{zz},\qquad E=H-F.
 \label{eq:ansatz-defect}
\end{equation}
The exact equation for the error is
\begin{equation}
 (\partial_t+\partial_{zz})E=-\mathfrak r_F.
 \label{eq:forced-error}
\end{equation}
Initial mismatch and evolution defect are separate data. In
particular, \(\mathfrak r_F=0\) does not identify the canonical
solution without a comparison of the initial fields. Conversely,
a field estimate need not control its time derivative merely because
its displayed allowance is differentiable. One can instead use
\begin{equation}
 E_t=-E_{zz}-\mathfrak r_F,
 \label{eq:direct-time-error}
\end{equation}
with independently proved spatial and residual estimates.

For a nonvanishing reference arm \(C\), put
\(p=C_z/C\), \(r_C=(C_t+C_{zz})/C\), and \(F=C(1+c)\).
The product rule gives exactly
\begin{equation}
 \frac{\mathfrak r_F}{C}
 =c_t+c_{zz}+2p c_z+r_C(1+c).
 \label{eq:correction-hierarchy}
\end{equation}
Cancelling successive orders in this expression generates a correction
hierarchy. Its integration data must match the gamma and arithmetic
initial expansion. The compact-positive correction in
Section~\ref{sec:cosine} is an existing example: its field theorem
uses the canonical integral as well as the residual calculation.

\subsection{Negative time gives forward diffusion}
Set \(T=-t\), \(h(T,z)=H(-T,z)\), \(f(T,z)=F(-T,z)\), and
\(e=h-f\). Then
\begin{equation}
 d=f_T-f_{zz}=-\mathfrak r_F(-T,z),\qquad
 e_T-e_{zz}=-d. \label{eq:negative-error}
\end{equation}
On the whole real line, in a function class for which the heat
convolution and differentiations are justified, define
\[
 (\mathsf S_v g)(x)=\frac1{\sqrt{4\pi v}}
 \int_\R e^{-(x-u)^2/(4v)}g(u)\dd u,\qquad v>0,
 \quad \mathsf S_0=I.
\]
Duhamel's formula \cite[Section~2.7]{duhamel} yields
\begin{equation}
 e(T)=\mathsf S_{T-T_0}e(T_0)
      -\int_{T_0}^{T}\mathsf S_{T-s}d(s)\dd s,
      \qquad T\ge T_0. \label{eq:duhamel}
\end{equation}
For bounded real-line data and integrable bounded forcing this implies
\begin{equation}
 \|e(T)\|_\infty\le\|e(T_0)\|_\infty
                  +\int_{T_0}^T\|d(s)\|_\infty\dd s.
 \label{eq:absolute-stability}
\end{equation}
This is an absolute estimate. Our field allowances are often relative
to an exponentially small carrier, and our formulas are defined on
buffered charts. Applying \eqref{eq:duhamel} at that stronger scale
requires a justified extension or localization, full Gaussian tails,
complex-strip losses, and a suitable weighted stability estimate.
For example a time-independent spatial cutoff \(\chi\) gives
\[
 (\partial_T-\partial_{xx})(\chi e)
 =-\chi d-2\chi'e_x-\chi''e.
\]
The two commutator terms cannot be omitted.

For a positive smooth weight \(\mathcal A(T,x)\), the normalized
error \(v=e/\mathcal A\) satisfies the exact real-line identity
\begin{equation}
 v_T-v_{xx}-2\frac{\mathcal A_x}{\mathcal A}v_x
 +\frac{\mathcal A_T-\mathcal A_{xx}}{\mathcal A}v
 =-\frac d{\mathcal A}. \label{eq:weighted-error}
\end{equation}
Its coefficients and boundary or tail data must be estimated in the
chosen region. No contraction in the outer-carrier norm is asserted.
Subtracting an unevaluated Duhamel defect from \(f\) simply recovers
the heat evolution of \(f(T_0)\); a useful asymptotic construction
must evaluate or approximate that correction with a quantitative gain.

Increasing \(T\) means moving toward more negative original time.
Neither \eqref{eq:absolute-stability} nor \eqref{eq:duhamel} supplies
stable recovery toward \(T=0\), or toward increasing positive \(t\)
for arbitrary initial errors. Positive-time analysis instead uses the
canonical analytic source, \eqref{eq:positive-anchor}, and controlled
contour/saddle comparisons. Thus the same residual bookkeeping applies
at all times, while the stability mechanism depends on the time region.

\subsection{Amplitude and phase as evolving quantities}
For a nonzero complex arm on a real interval, write
\(f=a e^{i\theta}\), with \(a>0\) and real unwrapped \(\theta\).
Then
\begin{equation}
 e^{-i\theta}(f_T-f_{xx})
 =a_T-a_{xx}+a\theta_x^2
  +i(a\theta_T-2a_x\theta_x-a\theta_{xx}).
 \label{eq:amplitude-phase-defect}
\end{equation}
Consequently an exact heat arm obeys
\begin{equation}
 a_T=a_{xx}-a\theta_x^2,\qquad
 \theta_T=\theta_{xx}+2(a_x/a)\theta_x.
 \label{eq:amplitude-phase-flow}
\end{equation}
For an approximate arm the real and imaginary parts of
\(e^{-i\theta}d\) are its forcing terms. Equivalently, for
\(f=f_0(1+c)\), \(p_0=(f_0)_x/f_0\) and
\(r_0=((f_0)_T-(f_0)_{xx})/f_0\),
\begin{equation}
 \frac{f_T-f_{xx}}{f_0}
 =c_T-c_{xx}-2p_0c_x+r_0(1+c).
 \label{eq:negative-correction-hierarchy}
\end{equation}
Solving the first correction equation with right side \(-r_0\)
leaves \(r_0c\), before further truncation. This is an exact algebraic
starting point, not a solvability or remainder result. Apply it to
named arms and retain their signed arithmetic sum and cutoff correction;
the full field has zeros and need not admit a global nonvanishing
polar representation.

\section{The outer carrier and its evolution correction}
\label{sec:cosine}
This elementary carrier isolates the first amplitude and phase.
It is a building block whose adequacy depends on the arithmetic
remainder at the chosen time, not a uniform all-time model. Let
\begin{equation}
\ell_X=\log\frac{X}{4\pi},\qquad
F_{\mathrm o}(t,X)=A_{\mathrm o}(t,X)\cos\Theta_{\mathrm o}(t,X),
\label{eq:outer}
\end{equation}
where
\begin{align}
A_{\mathrm o}(t,X)&=\frac{\pi^{1/4}}{16}X^{7/4}
\exp\left\{-\frac{\pi X}{8}
             +\frac t{16}\left(\ell_X^2-\frac{\pi^2}{4}\right)\right\},
\label{eq:amplitude}\\
\Theta_{\mathrm o}(t,X)&=\frac X4(\ell_X-1)
                         +\frac{\pi t}{16}\ell_X+\frac{7\pi}{8}.
\label{eq:phase}
\end{align}
The amplitude in \eqref{eq:amplitude} is the full cosine amplitude.
Each individual exponential carrier has half this amplitude.

\subsection{Calibration against the linear gamma model}
Define the exact logarithmic derivative
\begin{equation}
\alpha(w)=\frac{\Mc'(w)}{\Mc(w)}
=\frac1w+\frac1{w-1}-\frac{\log\pi}{2}+\frac12\psi(w/2),
\end{equation}
and the parent linear model
\begin{equation}
\Gc_t(w)=\Mc(w)e^{t\alpha(w)^2/4},\qquad
P_{\rm lin}(t,z)=\Gc_t(s)+\Gc_t(1-s). \label{eq:linear}
\end{equation}
Sectorial Stirling expansions \cite{dlmf-gamma}, at \(s_X=(1+iX)/2\),
give the following leading quantities:
\begin{align}
|\Mc(s_X)|&=\frac{\pi^{1/4}}{32}X^{7/4}e^{-\pi X/8}
                          (1+O(X^{-1})), \label{eq:gamma-amplitude}\\
\arg\Mc(s_X)&=\frac X4(\ell_X-1)+\frac{7\pi}{8}+O(X^{-1}),\\
\alpha(s_X)&=\frac{\ell_X}{2}+\frac{i\pi}{4}+O(X^{-1}).
\label{eq:alpha-leading}
\end{align}
The phase is an unwrapped, conjugate-compatible branch.
Indeed \(|s_X(s_X-1)|=(1+X^2)/4\); multiplying this by
\(\pi^{-1/4}|\Gamma(1/4+iX/4)|/16\) produces the constant \(1/32\).
The polynomial factor contributes the phase \(\pi\), while the gamma
factor contributes the constant \(-\pi/8\).
Finally
\[
\frac t4\left(\frac{\ell_X}{2}+\frac{i\pi}{4}\right)^2
=\frac t{16}\left(\ell_X^2-\frac{\pi^2}{4}\right)
 +i\frac{\pi t}{16}\ell_X.
\]
Pairing the conjugate carriers yields
\eqref{eq:amplitude}--\eqref{eq:phase}.

\begin{proposition}[Leading reduction on a fixed positive interval]
For a fixed \(32\le t\le T_0<\infty\), the parent linear-carrier
comparison and the preceding Stirling reduction give, on the real axis,
\begin{equation}
H(t,X)=F_{\mathrm o}(t,X)
 +A_{\mathrm o}(t,X)\,
        O_{T_0}\!\left(\frac{1+\log^2X}{X}\right).
\label{eq:outer-fixed}
\end{equation}
This statement retains the in-manuscript status of the imported
comparison. It asserts no relative error divided by the cosine at a zero.
\end{proposition}
\begin{proof}
For bounded \(t\), \eqref{eq:alpha-leading} changes the exponent in
\eqref{eq:linear} by \(O_{T_0}((1+\log X)/X)\).
The gamma amplitude and phase changes cost \(O(X^{-1})\).
These estimates compare each nonzero carrier to its elementary
counterpart. Add the two errors and then the parent error
\(O_{T_0}((1+\log^2X)/X)\) times the carrier sum.
\end{proof}

The compact complex-buffer, differentiated and first-correction version
on \(32\le t\le33\) is proved below at its in-note status. Constants
and onset gates are finite but not numerically instantiated. Expansion
and joining of the time ranges remain separate tasks.
The analytical floor \(32\) is not a physical transition. The published
small-positive-time asymptotics in \cite[Prop.~9.1]{polymath} provide a
separate starting range \(0<t\le1/2\), with an onset growing as
\(\exp(C/t)\); their hypotheses and errors must be reconciled before
joining the ranges. The interval between these imported ranges is not
declared covered by this draft.

\subsection{What a higher-order expansion must include}
A two-term formula is not obtained by expanding only the gamma factor.
The heat convolution, displacement of the saddle, and Gaussian
determinant contribute at the same orders. As a formal fixed-\(t\)
checkpoint, Gaussian completion suggests the correction
\begin{equation}
\log(\hbox{one gamma heat carrier})
\sim \log\Mc(w)+\frac t4\alpha(w)^2
       +\frac t4\alpha'(w)+\frac{t^2}{8}\alpha(w)^2\alpha'(w)+\cdots .
\label{eq:formal-correction}
\end{equation}
Ansatz Note03 proves this checkpoint, with a remainder, on
\(32\le t\le33\), as recorded below. Outside that bounded result it
remains a formal guide. Arithmetic modes may exceed a \(1/X\)
correction at small positive time. The next term must be selected by comparing all error
sources, rather than by automatically writing an inverse-height series.
At \(t=64\), \eqref{eq:phase} agrees with the leading full-saddle phase
in parent Note60, including \(4\pi\ell_X+7\pi/8\); agreement of those
leading coefficients does not identify the higher corrections.

\subsection{First corrected carrier on a compact positive interval}
Ansatz Note03 \cite{ansatz03} proves the preceding checkpoint on
\(32\le t\le33\), including the complete arithmetic and contour errors.
Fix \(\ell,h>0\). The spatial target includes
\( |\Re z-X|\le\ell+2\), \( |\Im z|\le1+h+2\), with one further
unit of analytic margin. On the right tail put
\[
 L=\Log(z/(4\pi)),\quad b_\pm=L/2\pm i\pi/4,\quad
 C_\pm=\tfrac12 A_{\mathrm o}(t,z)e^{\pm i\Theta_{\mathrm o}(t,z)},
 \qquad \mathcal E_o=|C_+|+|C_-|.
\]
The powers use the branch continued from positive real height. Define
\begin{equation}
 D_\pm=\frac{\pm i}{z}\left(
 \frac1{24}-\frac t4-\frac{7t b_\pm}{4}-\frac{t^2b_\pm^2}{8}\right),
 \qquad F_1=\sum_\pm C_\pm(1+D_\pm).
 \label{eq:corrected-outer}
\end{equation}
For sufficiently large \(X\), uniformly on the buffered domain,
\begin{align}
 |\partial_z^j(H-F_1)|&\le C_j\mathcal E_o
       \frac{1+\log^{4+j}X}{X^2},\qquad 0\le j\le3,
       \label{eq:corrected-jets}\\
 |\partial_t(H-F_1)|&\le C\mathcal E_o
       \frac{1+\log^6X}{X^2}.
       \label{eq:corrected-time}
\end{align}
Constants depend on the fixed buffers and are not numerically
instantiated. These are proved-in-note estimates with same-model
internal assessment, not independent specialist validation.

On the real axis the same field can be written, within the error in
\eqref{eq:corrected-jets}, as
\begin{align}
 H(t,X)&=A_{\mathrm o}e^{a_1/X}
              \cos(\Theta_{\mathrm o}+p_1/X)
              +A_{\mathrm o}O((1+\ell_X^4)/X^2),\\
 a_1&=\frac{7\pi t}{16}+\frac{\pi t^2\ell_X}{32},\\
 p_1&=\frac1{24}-\frac t4-\frac{7t\ell_X}{8}
       -\frac{t^2\ell_X^2}{32}+\frac{\pi^2t^2}{128}.
\end{align}
The derivation uses the actual fixed-endpoint gamma Gaussian integral.
Its logarithm equals the four terms in
\eqref{eq:formal-correction}, with remainder
\(O((1+\log^4X)/X^2)\). The determinant contributes
\(t\alpha'/4\); the displaced Gaussian mean contributes
\(t^2\alpha^2\alpha'/8\). On this interval the full arithmetic
remainder is \(O(X^{-t\log2/4})\), below the retained error.
The bounded middle and both infinite tails are paid separately.

For a direct consistency check put
\(Q(b)=21/16+t/8+5tb/8+t^2b^2/16\). The exact elementary residual is
\begin{equation}
 \frac{(\partial_t+\partial_{zz})C_\pm}{C_\pm}
 =\frac{\pm i(1+7b_\pm+tb_\pm^2)}{4z}+\frac{Q(b_\pm)}{z^2}.
 \label{eq:elementary-residual}
\end{equation}
The correction cancels its \(1/z\) term, leaving residual
\(O(\mathcal E_o(1+\log^4X)/X^2)\). The time comparison follows
directly from the heat equation, spatial jets and this actual residual;
it is not a differentiated big-O allowance. No physical covariance
transfer is inferred here.


\section{Canonical arithmetic data and the crossover through zero}
\label{sec:arithmetic}
\subsection{An exact normalization at zero time}
Let
\begin{align}
N_0(X)&=\left\lfloor\sqrt{X/(4\pi)}\right\rfloor,\\
\theta_\Gamma(X)&=\arg\Gamma(1/4+iX/4)-\frac X4\log\pi+\pi,\\
A_\Gamma(X)&=2|\Mc(s_X)|.
\end{align}
The classical Riemann--Siegel formula \cite[25.10.1--25.10.3]{dlmf-rs}
then gives
\begin{equation}
H(0,X)=A_\Gamma(X)
 \sum_{n=1}^{N_0(X)}\frac{\cos(\theta_\Gamma(X)-\tfrac X2\log n)}{\sqrt n}
 +A_\Gamma(X)\,O(X^{-1/4}).
\label{eq:RS-zero}
\end{equation}
To check the sign, write \(Z(r)=e^{i\vartheta(r)}\zeta(1/2+ir)\).
Since \(s_X(s_X-1)<0\), \(H(0,X)=-|\Mc(s_X)|Z(X/2)\);
adding \(\pi\) to \(\vartheta(X/2)\) gives \eqref{eq:RS-zero}.

The terms with \(n\ge2\) are part of the leading arithmetic expression.
No estimate in this draft licenses replacing their sum by zero.
Thus \eqref{eq:outer-fixed} is not extended uniformly to \(t=0\).

\subsection{A finite-sum ansatz near zero time}
A natural candidate retains these modes with heat-dependent weights:
\begin{align}
w_n(t,X)&=n^{-1/2}
 \exp\left\{\frac t4\big((\log n)^2-\ell_X\log n\big)\right\},
\label{eq:weights}\\
\varphi_n(t,X)&=\Theta_{\mathrm o}(t,X)
                   -\left(\frac X2+\frac{\pi t}{8}\right)\log n,\\
F_{\rm arith}(t,X)&=
 A_{\mathrm o}(t,X)\sum_{n\le N(t,X)}
                              w_n(t,X)\cos\varphi_n(t,X).
\label{eq:arithmetic-candidate}
\end{align}
These expressions are a proposed leading reduction. They do not include
an automatically small cutoff correction.

Here is the algebra behind the weights. Linearizing one gamma arm in
\eqref{eq:positive-anchor}, a Dirichlet term \(n^{-w-v}\) has Gaussian
factor
\[
\frac1{\sqrt{\pi t}}\int_\R
 e^{-v^2/t+v(\alpha(w)-\log n)}\dd v
 =\exp\left\{\frac t4(\alpha(w)-\log n)^2\right\}.
\]
Factoring out \(e^{t\alpha(w)^2/4}\) and using
\eqref{eq:alpha-leading} gives \eqref{eq:weights} and the displayed
phase shift. Interchanging this local calculation with a complete
Dirichlet sum or contour deformation requires its own justification.

Polymath's finite-sum construction \cite[Eqs.~(14)--(19)]{polymath}
uses \(N=\lfloor\sqrt{X/(4\pi)+t/16}\rfloor\), for
\(0<t\le1/2\), \(0\le y\le1\), \(X\ge200\). Its corrected version
retains a separate Riemann--Siegel term. That domain is a positive-time
input, not a theorem for \(t<0\). For holomorphic derivatives, use a
cutoff fixed on each analytic chart, or a proved corrected
representation across cutoff changes. Differentiating the floor
function as though the resulting field were holomorphic is invalid.

\subsection{A candidate time scale, not a universal transition law}
For fixed \(n\), set \(\tau=t\ell_X\). If \(t=\tau/\ell_X\), then
\begin{equation}
w_n(t,X)\longrightarrow n^{-1/2-\tau/4}. \label{eq:tau}
\end{equation}
This makes \(t\log X\) a natural coordinate for the loss of fixed
arithmetic modes. It does not prove a one-variable limiting profile:
\(N\) grows with \(X\), and the collective contribution near the cutoff
must also be estimated. No finite universal critical value of \(\tau\)
is asserted. At bounded \(\tau\), even the fixed \(n=2\) coefficient
does not tend to zero. Negative \(\tau\) increases its relative weight.

The first crossover task is therefore a uniform finite-sum comparison
through zero time, not a fitted sigmoid between two cosines.

\subsection{A retained correction and a signed real-axis comparison}
The following bounded result is proved in Ansatz Note03
\cite{ansatz03}. It does not supply the full buffered complex theorem.
On a fixed far-height chart choose
\(N=\lfloor\sqrt{X/(4\pi)}\rfloor\) once, and set
\begin{align}
 a(z)&=\sqrt{z/(4\pi)},&
 G(z)&=\sqrt{\Mc(s(z))\Mc(1-s(z))},\\
 \mathfrak C(p)&=
 \frac{\cos(2\pi(p^2-p)-\pi/8)}{\cos(2\pi p)},&
 \mathcal R_N(z)&=(-1)^NG(z)a(z)^{-1/2}
                           \mathfrak C(a(z)-N).
 \label{eq:fixed-rs-correction}
\end{align}
Choose \(G\) positive on the real segment. All apparent poles of
\(\mathfrak C\) are removable. Its recurrence
\[
 \mathfrak C(p)+\mathfrak C(p-1)
       =2\cos(2\pi(p-1)^2-\pi/8)
\]
pays the adjacent cutoff mode. At zero time this yields a corrected
real-axis remainder \(O(GX^{-3/4})\), uniformly across a cutoff.

Use the weights and phases in \eqref{eq:weights} and define
\begin{equation}
 F_N(t,u)=A_{\mathrm o}(t,u)\sum_{n\le N}w_n(t,u)\cos\varphi_n(t,u)
                    +e^{-\pi^2t/64}\mathcal R_N(u).
 \label{eq:common-real-bridge}
\end{equation}
For real \(u\) in a fixed horizontal target window and sufficiently
large chart center, uniformly for \(|t|\log(u/(4\pi))\le8\),
\begin{equation}
 |H(t,u)-F_N(t,u)|\le C A_{\mathrm o}(t,u)u^{-1/4}.
 \label{eq:real-bridge-bound}
\end{equation}
Finite analytical constants and onset gates are specified in Note03;
no numerical onset is certified.

For \(t=-T<0\), the proof uses the exact real Gaussian convolution,
the full \(|v|\le16\) window, gamma-envelope tails up to
\(|v|=u/2\), and the complete-source bound beyond it. All modes
through \(N\) are paid using
\(\sum w_n\le C u^{1/4}e^{T\ell_u^2/16}\).
The corrected source remainder costs at most
\(u^{-3/4}e^{T\ell_u^2/16}\le C u^{-1/4}\) at the requested edge.
The bare \(u^{-1/4}\) Riemann--Siegel error would not give decay
throughout this window. Smoothing the correction gives the displayed
factor \(e^{\pi^2T/64}\), within the proved error.

For \(0<t\ell_u\le8\), Polymath Theorem~1.3 in its real-axis
domain pays the field, gamma replacement and adjacent cutoff term.
Adding the explicit correction in \eqref{eq:common-real-bridge} is
within the same coarse error. This does not identify it with the full
positive-time \(C_t\) at finer precision.

\subsection{Complex remainder and the full buffered crossover}
Ansatz Note04 \cite{ansatz04} closes the preceding complex estimate
without an \(X^\epsilon\) loss. For any fixed radii \(r,b>0\),
\begin{equation}
 |H(0,z)-S_N(z)-\mathcal R_N(z)|
 \le C_{r,b}X^{-3/4}\{ |\Mc(s)|+|\Mc(1-s)|\},
 \quad |\Re z-X|\le r,\quad |\Im z|\le b.
 \label{eq:buffer-rs}
\end{equation}
Here \(N=\lfloor\sqrt{X/(4\pi)}\rfloor\) remains fixed on the chart.
Constants and finite onset gates depend on the declared fixed buffers;
they are not numerically instantiated. The result is proved in-note
from named classical inputs, with same-model internal assessment only.

The proof pays the analytic boundary rather than extending real values
by assertion. Set \(A=\log^2X\), \(D=\max(4,b+1)\), and normalize
the error by the nonzero dominant arm \(\Mc(1-s)\) on the rectangle
\( |\Re z-X|\le A\), \(0\le\Im z\le D\). The corrected real
Riemann--Siegel formula and adjacent-cutoff recurrence give the lower
edge bound \(O(X^{-3/4})\). On the upper edge, combine
\(H/\Mc(1-s)\) with its dominant finite sum as one ordinary Dirichlet
tail. That tail and the reflected finite sum are
\(O(X^{(1-D)/4})\), hence \(O(X^{-3/4})\); the correction is smaller.
Euler summation bounds the vertical edges by \(O(X^{1/2})\).
The analytic multiplier \(X^{3/4}e^{-(z-X)^2}\) suppresses those
edges. The maximum principle and reflection prove
\eqref{eq:buffer-rs} on the fixed target core. The growing real-edge
cutoff error and all onset gates are given in Note04.

Use the same \(F_N\) as in \eqref{eq:common-real-bridge}, now
holomorphically defined on a fixed chart. Let
\[
 L_X=\log(X/(4\pi)),\quad \Lambda_X=1+L_X,\quad
 t_-=\max(-t,0),\qquad
 \delta(t,X)=X^{-3/4}e^{t_-L_X^2/16}.
\]
For the full fixed physical, averaging and analytic buffers, uniformly
for \(|t|L_X\le8\) and sufficiently large \(X\),
\begin{align}
 |\partial_z^j(H-F_N)|&\le C_j\mathcal E_o\delta\Lambda_X^j,
                       \qquad 0\le j\le3,
                       \label{eq:buffered-crossover-jets}\\
 |\partial_t(H-F_N)|&\le C\mathcal E_o\delta\Lambda_X^2.
                       \label{eq:buffered-crossover-time}
\end{align}
In particular \(\delta\le C X^{-1/4}\) throughout the signed window,
and \(\delta=X^{-3/4}\) at nonnegative times, including zero.
The physical choice is \(r=\ell+2\), \(b=1+h+2\), with an
additional analytic margin. No cosine is used as a denominator.

For negative time the exact real Gaussian convolution consumes
\eqref{eq:buffer-rs}; for positive time the exact imaginary Gaussian
convolution does so. Both use a fixed window \(V_b\), chosen sufficiently
large from the declared source/ray constants, with initial radius
\(r+V_b+2\) or strip \(b+1+2V_b\), respectively. Near tails retain
the gamma height exponential and distant tails are separately paid.
All \(N\) modes are counted through
\[
 \sum_{\pm,n\le N}|P_{\pm,n}(t,z)|
 \le C\mathcal E_o X^{1/4}e^{t_-L_X^2/16},
\]
where \(P_{\pm,n}\) are the two elementary mode arms.
This proves the field error and \(1/\log X\) Cauchy disks prove the
spatial jets. The positive-time proof does not extend the domain of
the published evolved remainder theorem by assertion.

The direct time estimate follows from an exact residual. With
\(c=b_\pm-\log n\),
\begin{equation}
 \frac{(\partial_t+\partial_{zz})P_{\pm,n}}{P_{\pm,n}}
 =\frac{\pm i(1+7c+tc^2)}{4z}
 +\frac{21/16+t/8+5tc/8+t^2c^2/16}{z^2}.
 \label{eq:mode-residual}
\end{equation}
The correction satisfies
\(\mathcal R_N''-(\pi^2/64)\mathcal R_N=O(|G|X^{-3/4})\).
Their sum is bounded by \(C\mathcal E_o\delta\Lambda_X\).
The identity
\[
 \partial_t(H-F_N)=-(H-F_N)_{zz}
                   -(\partial_t+\partial_{zz})F_N
\]
then proves \eqref{eq:buffered-crossover-time} through zero, without a
\(1/|t|\) loss or differentiation of an error allowance.

This closes the bounded AN01--AN02 field/jet/time request. A globally
floor-valued formula is not silently differentiated; overlapping fixed
charts compare through the same \(H\). Physical transfer and join
defects, independent specialist review, and all-time assembly remain
separate. In particular the factor \(e^{TL_X^2/16}\) prevents this
fixed-order budget from covering a fixed compact interval of negative
times by substitution.


\section{Negative-time evolution and the two accuracy scales}
\label{sec:negative}
For fixed \(T>0\), the source identity gives the exact real-initial-value
formula, also for complex \(z\),
\begin{equation}
 H(-T,z)=\frac1{\sqrt{4\pi T}}\int_{\R}
 e^{-(z-u)^2/(4T)}H(0,u)\dd u.
 \label{eq:negative-real-initial}
\end{equation}
Absolute Fubini follows from the real Fourier source and the Gaussian.
On \(|\Im z|\le b\) the kernel modulus costs \(e^{b^2/(4T)}\).
This is bounded on \(1\le T\le2\), the fixed interval in
Note05. The later bridge explicitly pays its small-\(T\) growth.
No contour deformation or unproved negative-time version of a
positive-time remainder theorem is used. A real part may remain inside
the real initial data; it is never taken after applying a complex kernel.

\subsection{A short formula at the natural gamma scale}
Ansatz Note05 \cite{ansatz05} extends the same short corrected formula
\(F_N\) of \eqref{eq:common-real-bridge} to fixed negative times at a
different comparison scale. Set
\[
 G(z)=\{\Mc(s)\Mc(1-s)\}^{1/2},\qquad
 L=\log(X/(4\pi)),\qquad N=\lfloor\sqrt{X/(4\pi)}\rfloor,
\]
using the branch positive on the real segment. For all sufficiently
large \(X\), uniformly on every full fixed buffered chart and
\(1\le T\le2\),
\begin{align}
 |\partial_z^j(H-F_N)(-T,z)|&\le C_j|G(z)|X^{-3/4},
                &&0\le j\le3,\label{eq:negative-natural}\\
 |\partial_t(H-F_N)(-T,z)|&\le C|G(z)|X^{-3/4}.
 \label{eq:negative-natural-time}
\end{align}
The physical chart has radii \(r=\ell+2\), \(b=1+h+2\), with an
extra fixed analytic margin. On such charts \(G'/G=-\pi/8+O(X^{-1})\),
so fixed-radius Cauchy disks incur no logarithmic jet loss.

The proof smooths the real corrected fixed-cutoff remainder on
\(|u-X|\le8L\), within Note04's growing real chart. Per-mode logarithmic
curvature is \(O(X^{-1})\), independently of \(n\). The transverse
mode powers cancel when comparison is made against real initial data.
All logarithmic moments of the elementary heat-mode population satisfy
\[
 \sum_{\pm,n\le N}(1+|b_\pm-\log n|)^d|P_{\pm,n}(-T,z)|
 \le C_{b,d}|G(z)|X^{1/4},\qquad d\ \hbox{fixed}.
\]
This pays both the full mode count and gamma-to-elementary replacement.
The correction smooths to \(e^{\pi^2T/64}\mathcal R_N\), within the
stated error. Its residual and \eqref{eq:mode-residual}, summed with
these moment bounds, prove the direct time estimate. Genuine and
artificial Gaussian tails are separately paid in Note05.

Relative to the positive correction envelope \(|G|X^{-1/4}\), the
error is \(O(X^{-1/2})\). This is not a lower bound on \(H\) or
\(\mathcal R_N\), which can vanish. At the real outer-carrier scale
the allowance becomes
\(X^{-3/4}e^{TL^2/16}\), which grows at fixed \(T\).
This is a scoped failure of that bound, not a lower bound on the error.

\subsection{Intermediate negative times as evolving arithmetic modes}
The short model also has an informative exact rewriting, derived in
Note08 \cite{ansatz08}. On the real right tail let
\[
 \begin{gathered}
 a_X=\sqrt{X/(4\pi)},\quad r_n=\log(a_X/n),\quad k_n=r_n/2,\\
 A_0=A_{\mathrm o}(0,X),\quad
 \phi_n(0,X)=\Theta_{\mathrm o}(0,X)-(X/2)\log n.
 \end{gathered}
\]
Completing the square in its weights gives
\begin{align}
 A_{\mathrm o}(-T,X)w_n(-T,X)
 &=A_0n^{-1/2}e^{\pi^2T/64}e^{-Tk_n^2},\\
 \phi_n(-T,X)&=\phi_n(0,X)-\pi T k_n/4.
 \label{eq:negative-mode-weights}
\end{align}
Hence the existing corrected short formula is exactly
\begin{equation}
 F_N(-T,X)=e^{\pi^2T/64}\left[
 A_0\sum_{n\le N}n^{-1/2}e^{-Tk_n^2}
 \cos\{\phi_n(0,X)-\pi T k_n/4\}
 +\mathcal R_N(X)\right].
 \label{eq:negative-mode-sum}
\end{equation}
This identity adds no error estimate or validity domain to
\eqref{eq:negative-natural}. For a fixed mode,
\(\partial_X\phi_n(0,X)=k_n\). A locally frozen exponential with
\(\lambda_n=-\pi/8+ik_n\) evolves by
\[
 e^{T\lambda_n^2}
 =e^{\pi^2T/64}e^{-Tk_n^2}e^{-i\pi T k_n/4}.
\]
This explains the displayed damping and phase drift. The slowly
varying amplitude and spatial phase curvature create a residual and
are natural inputs to \eqref{eq:negative-correction-hierarchy}.

There are several distinct scales. For fixed \(n\), the damping
factor changes when \(T\log^2X=O(1)\). Ratios of two fixed-mode
weights change when \(T\log X=O(1)\). The damping band
\(r_n=O(T^{-1/2})\), however, contains a growing population with
\(n\) comparable to \(\sqrt X\) at fixed \(T\).
These are statements about components. They do not identify a dominant
term of the signed sum, justify dropping the RS correction, or establish
a universal transition time. The full-population bounds used in
Note05 must remain in any proposed correction.

\subsection{One evaluated correction at fixed order}
\label{sec:negative-first-correction}
Ansatz Note10, assessed in Review07 \cite{ansatz10,review07}, now pays
the first natural-scale improvement. Freeze \(L,N\) as above and
\(V=8L\), \(I_X=[X-V,X+V]\). The chart radii are
\(r=\ell+2\), \(b=1+h+2\), with an extra analytic unit.
For \(\epsilon=\pm1\), \(n\le N\), define target-dependent exact modes
\begin{align}
 w_\epsilon(z)&=(1+\epsilon iz)/2,&
 h_{\epsilon n}(z)&=\Mc(w_\epsilon(z))n^{-w_\epsilon(z)},\\
 \beta_{\epsilon n}(z)&=\frac{\epsilon i}{2}
             [\alpha(w_\epsilon(z))-\log n],&
 \gamma_\epsilon(z)&=\beta_{\epsilon n}'(z)
                    =-\alpha'(w_\epsilon(z))/4.
 \label{eq:first-mode-data}
\end{align}
The curvature \(\gamma_\epsilon\) is independent of \(n\).
The evaluated first evolution correction is
\begin{equation}
 \mathcal A_{\epsilon n}(T,z)
 =h_{\epsilon n}(z)e^{T\beta_{\epsilon n}(z)^2}
 \left[1+\gamma_\epsilon(z)
          \{T+2T^2\beta_{\epsilon n}(z)^2\}\right].
 \label{eq:first-evaluated-correction}
\end{equation}
It matches the exact gamma mode at \(T=0\).

The initial RS datum must also improve. For real \(u\in I_X\) put
\(a_u=\sqrt{u/(4\pi)}\), \(m=m(u)=\lfloor a_u\rfloor\),
\(p_u=1-2(a_u-m)\), \(s_u=(1+iu)/2\), and
\(U_u=e^{-i\{(u/4)\log(u/(4\pi))-u/4-\pi/8\}}\).
Retain \(K=1\) in Polymath Proposition~6.2 \cite{polymath}:
\begin{align}
 R_{m,1}(u)&=2\Re\left\{(-1)^{m-1}\Mc(s_u)U_u a_u^{-1/2}
              [C_0(p_u)+a_u^{-1}C_1(p_u,1/2)]\right\},\\
 Q_{X,1}(u)&=S_{m(u)}(u)-S_N(u)+R_{m(u),1}(u),\\
 \mathcal I_1(T,z)&=\frac1{\sqrt{4\pi T}}\int_{I_X}
                 e^{-(z-u)^2/(4T)}Q_{X,1}(u)\dd u.
 \label{eq:first-rs-integral}
\end{align}
The \emph{exact} adjacent transfer \(S_m-S_N\) remains in the real
initial data. Under \(a(X+V)-a(X-V)<1\), at most one threshold occurs.
Its partitions and both endpoints are frozen independently of \(z,T\),
so \((\mathcal I_1)_T=(\mathcal I_1)_{zz}\) exactly. No regularity or
endpoint identity for \(C_1\) is inferred from its size bound.
Define the fixed-order representation
\begin{equation}
 F_X^{[1]}(-T,z)=\sum_{\epsilon=\pm1}\sum_{n\le N}
                  \mathcal A_{\epsilon n}(T,z)+\mathcal I_1(T,z).
 \label{eq:first-negative-model}
\end{equation}
This is a Gaussian-moment formula with a retained bounded RS integral.
It contains all \(N\) modes. The integral uses named initial
coefficients and exact transfers, not an unknown remainder or a full
unevaluated Duhamel forcing.

\begin{proposition}[Fixed-order natural-scale improvement]
\label{prop:first-negative-correction}
Uniformly for \(1\le T\le2\), on the full fixed buffered charts, for
sufficiently large \(X\) depending on those buffers,
\begin{align}
 |\partial_z^j(H-F_X^{[1]})(-T,z)|
       &\le C_j|G(z)|X^{-5/4},&&0\le j\le3,
       \label{eq:first-negative-jets}\\
 |\partial_t(H-F_X^{[1]})(-T,z)|
       &\le C|G(z)|X^{-5/4},
       \label{eq:first-negative-time}\\
 |\mathfrak r_{F_X^{[1]}}|+|\partial_z\mathfrak r_{F_X^{[1]}}|
       &\le C|G(z)|X^{-7/4},
       \label{eq:first-negative-residual}
\end{align}
where \(\mathfrak r_F=F_t+F_{zz}\). There is no logarithmic loss.
Status is proved in-note from named classical inputs, with same-model
internal assessment only. Constants and onset are not numerically
instantiated. No RH or simplicity premise is used.
\end{proposition}

\begin{proof}
First derive the correction, separately from proving solution error.
Suppress the mode indices; set \(f_0=he^{T\beta^2}\),
\(\gamma=\beta'\), \(\eta=\beta''\), and \(p=(f_0)_z/f_0\).
Direct differentiation gives
\begin{align}
 p&=\beta+2T\beta\gamma,\\
 r_0:=\frac{(f_0)_T-(f_0)_{zz}}{f_0}
   &=-\gamma(1+4T\beta^2)+r_2,\\
 r_2&=-2T(\gamma^2+\beta\eta)-4T^2\beta^2\gamma^2.
 \label{eq:first-product-data}
\end{align}
The product equation \eqref{eq:negative-correction-hierarchy} selects
\(c=\gamma(T+2T^2\beta^2)\), since
\(c_T=\gamma(1+4T\beta^2)\). The full residual is exactly
\begin{equation}
 \frac{\mathcal A_T-\mathcal A_{zz}}{f_0}
       =r_2+r_0c-c_{zz}-2pc_z.
 \label{eq:first-exact-residual}
\end{equation}
In particular transport and coefficient derivatives have not vanished.
Sectorial gamma bounds give
\(\gamma=O(X^{-1})\), \(\eta=O(X^{-2})\),
\(\beta'''=O(X^{-3})\). More explicitly,
\[
 \begin{aligned}
 A&=T+2T^2\beta^2,&c_z&=\eta A+4T^2\beta\gamma^2,\\
 c_{zz}&=\beta'''A+12T^2\beta\gamma\eta+4T^2\gamma^3.
 \end{aligned}
\]
Thus \(|\mathcal A_T-\mathcal A_{zz}|
\le C|f_0|X^{-2}(1+|\beta|^4)\).

For every fixed integer \(d\ge0\), the \emph{complete} population obeys
\begin{equation}
 \sum_{\epsilon,n\le N}|h_{\epsilon n}e^{T\beta_{\epsilon n}^2}|
       (1+|\beta_{\epsilon n}|)^d
       \le C_{b,d}G(X)X^{1/4}.
 \label{eq:first-population}
\end{equation}
Indeed \(\Re\beta=-\pi/8+O(X^{-1})\) and
\(\Im\beta=\epsilon r_n/2+O_b(X^{-1})\).
Off-axis powers combine into \(e^{C_b r_n}\); a mode is bounded by
\(CG(X)a_X^{-1/2}e^{C_b r_n-r_n^2/8}\).
The shell \(j\le r_n<j+1\) contains at most \(2a_Xe^{-j}\)
integers. Summing every fixed polynomial moment proves
\eqref{eq:first-population}. It pays the residual just computed.
Fixed-margin Cauchy on the summed holomorphic residual gives its first
spatial derivative, proving \eqref{eq:first-negative-residual}; original
time residual is the negative of the \(T\)-heat defect.

For the solution comparison, the classical \(K=1\) initial remainder
is independently
\begin{equation}
 |H(0,u)-S_{m(u)}(u)-R_{m(u),1}(u)|
 \le \frac{2^{7/4}1.1^2}{7}G(u)a_u^{-5/2}
 \le CG(u)X^{-5/4}.
 \label{eq:first-initial-remainder}
\end{equation}
Split \eqref{eq:negative-real-initial} at \(I_X\) and insert this
datum with \(S_N+Q_{X,1}\). On the enlarged target chart the complex
kernel modulus contributes bounded \(e^{(b+1)^2/(4T)}\), and moments
against \(G(u)/G(X)\le Ce^{\pi|u-X|/8}\) are bounded on
\([1,2]\). The integrated initial error is therefore
\(CG(X)X^{-5/4}\).

For each arithmetic mode let \(v=u-z\). Sectorial Taylor expansion gives
\begin{equation}
 \frac{h(u)}{h(z)e^{\beta(z)v}}
 =1+\frac{\gamma(z)v^2}{2}
 +O\!\left(\frac{|v|^3+|v|^4}{X^2}e^{C|v|^2/X}\right).
 \label{eq:first-local-remainder}
\end{equation}
For \(z=x+iy\), comparison to the real initial line retains the
transverse cancellation
\[
 |h(z)e^{-iy\beta(z)}|\le C_b|h(x)|=C_bG(x)n^{-1/2},
 \qquad |\Re\beta(z)|\le C_b.
\]
Bounded weighted Gaussian moments give arithmetic error
\(CG(X)n^{-1/2}X^{-2}\) per mode. The entire initial population
\(\sum_{n\le N}n^{-1/2}=O(X^{1/4})\) costs
\(CG(X)X^{-7/4}\). Completing only the tangent-quadratic polynomial
to the real line evaluates its quadratic moment as
\((2T+4T^2\beta^2)e^{T\beta^2}\), exactly
\eqref{eq:first-evaluated-correction}.

All completions have paid tails. On
\(V<|u-X|\le X/2\), retain the true initial bound
\(CX^P e^{-\pi X/8+\pi|u-X|/8}\), for a fixed \(P\).
Gaussian dominance gives at most
\(CG(X)X^{P-7/4}(1+V^5)e^{-V^2/64}\).
Only beyond \(X/2\) use the global complete-source bound, where the
relative cost is \(CX^C e^{-cX^2+\pi X/8}\).
Artificial tangent-quadratic tails cost at most
\(CG(X)X^{1/4}(1+V^7)e^{-V^2/64}\), including all modes.
Since \(V=8L\), every tail is eventually below the required inverse
power. No true gamma mode is extended through remote poles, and no
analytic RS correction is continued along an unbounded ray.
Together with the initial remainder this proves the field estimate on
the enlarged chart. Note10 gives the explicit eventual smallness gates.

The target-centered polynomial depends on \(z\); spatial derivatives
do not merely hit its kernel. Instead fixed-radius Cauchy on the
holomorphic error, with \(|G|\) varying by only a bounded factor,
proves \eqref{eq:first-negative-jets} through three without log loss.
At fixed \(z\), direct \(T\) differentiation does hit only the
kernel and inserts the bounded-moment polynomial
\((z-u)^2/(4T^2)-1/(2T)\). Alternatively, the already proved
second spatial error and residual give
\(\partial_t(H-F)=-(H-F)_{zz}-\mathfrak r_F\), proving
\eqref{eq:first-negative-time}. No residual estimate was substituted
for a solution comparison, and no asymptotic allowance was differentiated.
The fixed initial-window split is not a spatially localized error PDE;
there is no omitted commutator or whole-line supremum contraction.
\end{proof}

An equivalent elementary-arm display, at the theorem's accuracy, retains
every \(P_{\epsilon n}\) and sets \(B=b_\epsilon-\log n\):
\begin{equation}
 D_{\epsilon n}(t,z)=\frac{\epsilon i}{z}
       \{1/24-t/4-(7t/4)B-(t^2/8)B^2\}.
 \label{eq:first-elementary-correction}
\end{equation}
Replacing \(\mathcal A_{\epsilon n}\) by
\(P_{\epsilon n}(-T,z)[1+D_{\epsilon n}(-T,z)]\) costs, in total,
\(C|G|X^{-7/4}\), including the required jets. It keeps both the
initial gamma phase and slope corrections. The old
\(e^{\pi^2T/64}\mathcal R_N\) is replaced by \(\mathcal I_1\).
The elementary product residual cancels its leading \(1/z\) term;
Note10 EC26 and the full moment bound give the same residual and direct
time allowances. Keeping only the old leading RS remainder would leave
\(C|G|X^{-3/4}\), irrespective of improved mode residuals.

The near-cutoff population still has order \(\sqrt X\) modes and
absolute arm weight comparable to \(GX^{1/4}\); first mode corrections
total \(O(GX^{-3/4})\). Neither is a leading asymptotic for the signed
sum. The integral in \eqref{eq:first-rs-integral} remains: an
integral-free short correction needs separately paid \(C_1\) regularity,
cutoff transfer and heat approximation, and no useful cost theorem follows.

Finally \(\mathcal E_o(-T,z)\ge c|G(z)|e^{-TL^2/16}\) converts the
solution and residual bounds only to the respective allowances
\begin{equation}
 \delta_1=X^{-5/4}e^{TL^2/16},\qquad
 \delta_{\mathrm{res},1}=X^{-7/4}e^{TL^2/16}.
 \label{eq:first-remaining-outer-loss}
\end{equation}
Their logarithms are \(TL^2/16-5L/4+O(1)\) and
\(TL^2/16-7L/4+O(1)\). Both grow at each fixed \(T\in[1,2]\).
The solution allowance improves by \(X^{-1/2}\), but no vanishing
outer-relative error is proved. This limits these finite algebraic
allowances, not all short resummations; it is no actual error lower bound.
The growing-order result below retains its stronger outer accuracy
and zero frozen-chart residual.

\subsection{Every declared fixed compact-negative-time family}
\label{sec:compact-negative-family}
Note12, assessed in Review09 \cite{ansatz12,review09}, extends the
exact-gamma correction to every declared pair of fixed constants
\(0<\tau\le B<\infty\). This is a uniform statement for
\(\tau\le T\le B\), not a substitution of moving endpoints into a
fixed-time estimate. Freeze
\begin{equation}
 \overline B=\max(1,B),\qquad V_B=8\sqrt{\overline B}\,L,
 \qquad I_{X,B}=[X-V_B,X+V_B].
 \label{eq:compact-negative-window}
\end{equation}
Use the same \(N\), exact-gamma modes
\eqref{eq:first-mode-data}, evaluated corrections
\eqref{eq:first-evaluated-correction}, and real RS/exact-transfer
datum \(Q_{X,1}\) in \eqref{eq:first-rs-integral}. Define
\begin{equation}
 F_{X,B}^{[1]}(-T,z)=\sum_{\epsilon=\pm1}\sum_{n\le N}
 \mathcal A_{\epsilon n}(T,z)
 +\int_{I_{X,B}} k_T(z-u)Q_{X,1}(u)\dd u,
 \quad k_T(v)=\frac{e^{-v^2/(4T)}}{\sqrt{4\pi T}}.
 \label{eq:compact-negative-model}
\end{equation}
All endpoints and real initial partitions are frozen in \(z,T\).
Every arithmetic mode, the \(K=1\) RS data and exact adjacent transfer
remain; the bounded integral again has zero heat residual. This family
is not identified with \(F_X^{[1]}\) when the windows differ.

\begin{proposition}[Endpoint-dependent compact-negative family]
\label{prop:compact-negative-family}
For declared fixed \(0<\tau\le B<\infty\), the sufficient onset
ledger in Appendix~\ref{app:compact-endpoints} yields, uniformly on
\(\tau\le T\le B\) and the full fixed target charts,
\begin{align}
 |\partial_z^j(H-F_{X,B}^{[1]})(-T,z)|
   &\le C_{r,b,j}\mathcal E\,|G(z)|X^{-5/4},&&0\le j\le3,\\
 |\partial_t(H-F_{X,B}^{[1]})(-T,z)|
   &\le C_{r,b}\mathcal E_t\,|G(z)|X^{-5/4},\\
 |\mathfrak r_{F_{X,B}^{[1]}}|
       +|\partial_z\mathfrak r_{F_{X,B}^{[1]}}|
   &\le C_{r,b}\mathcal R\,|G(z)|X^{-7/4}.
 \label{eq:compact-negative-bounds}
\end{align}
The explicit endpoint factors \(\mathcal E,\mathcal E_t,\mathcal R\)
are defined in \eqref{eq:compact-error-factors}; all remaining constants
depend only on the fixed buffers and named classical inputs, not on
\(\tau,B\). There is no logarithmic loss or numerically evaluated
onset. Status is proved in-note from the named inputs of Notes05/10,
with same-model internal analytical assessment only.
\end{proposition}

\begin{proof}
The exact canonical convolution \eqref{eq:negative-real-initial}
is split at \(I_{X,B}\). On that window the independent initial
\(K=1\) remainder \eqref{eq:first-initial-remainder} contributes
\(CG(X)X^{-5/4}\mathcal M_0\); its direct time derivative costs
\(CG(X)X^{-5/4}\mathcal T_0\).
For each arithmetic mode use \eqref{eq:first-local-remainder} and
the curvature gate in \eqref{eq:compact-local-onsets}, which bounds
its remainder exponential by \(e\). Transverse cancellation
\(|h(z)e^{-iy\beta(z)}|\le C|h(x)|=CG(x)n^{-1/2}\) precedes
absolute estimation. The local arithmetic comparison costs
\(CG(X)n^{-1/2}X^{-2}(\mathcal M_3+\mathcal M_4)\) per mode;
summing every \(n\le N\) gives \(GX^{-7/4}\).
Direct time uses \(\mathcal T_3+\mathcal T_4\), including moments
through six. The evaluated quadratic Gaussian moment is unchanged.

Appendix~\ref{app:compact-endpoints} pays the full evolved mode
population and both true and artificial tails before any final
allowance. The genuine intermediate tail retains its gamma height
exponential; only beyond distance \(X/2\) is global source
boundedness used. The two initial-window complements are paid, not
discarded as localization commutators. Fixed-margin Cauchy on the
holomorphic error gives spatial jets through three without log loss.
The target-centered polynomial depends on \(z\), so those derivatives
cannot be assigned to the kernel alone.

The exact product identity \eqref{eq:first-exact-residual}, with
the full population \eqref{eq:compact-population}, gives the last
bound; the integral contributes zero. Cauchy on the summed residual
gives its first spatial derivative. The independently proved second
spatial error and \(\partial_t(H-F)=-(H-F)_{zz}-\mathfrak r_F\)
also yield a time factor \(C_{r,b}(\mathcal E+\mathcal R)\).
No residual bound replaces the canonical solution comparison, and
no error allowance is differentiated.
\end{proof}

The carrier is re-derived on this interval. The gate
\(C_0B(1+L)/X\le1\), rather than an imported \(T\le2\) constant,
gives
\begin{equation}
 \mathcal E_o(-T,z)\ge c_{r,b}|G(z)|
       e^{-TL^2/16+\pi^2T/64}.
 \label{eq:compact-carrier-conversion}
\end{equation}
The solution and residual outer-relative allowances therefore retain
\begin{equation}
 X^{-5/4}e^{TL^2/16-\pi^2T/64},\qquad
 X^{-7/4}e^{TL^2/16-\pi^2T/64},
 \label{eq:compact-outer-loss}
\end{equation}
multiplied by the respective paid endpoint factors. They still grow
at every fixed \(T>0\), limiting these finite algebraic bounds without
giving an actual error lower bound or a universal no-go for resummation.
The factors in the appendix deteriorate as \(\tau\downarrow0\) or
\(B\uparrow\infty\); no endpoint-independent onset, growing-time
substitution or new through-zero join follows.

On \([\tau,B]\cap[1,2]\), after \emph{both} height onsets, compare
\(F_{X,B}^{[1]}\) and \(F_X^{[1]}\) through the same canonical
\(H\), with the sum of their proved allowances. Different windows
are not identified, even though their integral difference is itself
an exact heat solution. The elementary display
\eqref{eq:first-elementary-correction}, the older growing-order model
and the time join below retain their existing separate scopes.
This theorem concerns the exact-gamma family, with a retained RS
integral and no useful-cost or physical-transfer claim.

\subsection{A finite representation at outer-carrier precision}
The stronger outer-scale target is paid by the earlier growing-order
construction \cite{ansatz05}, which remains distinct from
\eqref{eq:first-negative-model}. It controls both the initial RS
remainder and arithmetic heat curvature. Freeze at the center
\begin{equation}
 V=8L,\qquad K=\lceil L\rceil,\qquad J=\lceil3L\rceil,
 \qquad R=X^{1/3},\qquad I_X=[X-V,X+V].
 \label{eq:negative-orders}
\end{equation}
For real \(u\in I_X\), put
\[
 a_u=\sqrt{u/(4\pi)},\quad m(u)=\lfloor a_u\rfloor,
 \quad p_u=1-2(a_u-m(u)),\quad s_u=\tfrac12+\tfrac{iu}2,
\]
\[
 U_u=\exp\{-i[(u/4)\log(u/(4\pi))-u/4-\pi/8]\}.
\]
Let \(C_k(p,1/2)\) denote the classical RS coefficients in
Polymath Proposition~6.2 \cite{polymath}. Define real initial data
\begin{align}
 R_{m,K}(u)&=2\Re\left\{(-1)^{m-1}\Mc(s_u)U_u a_u^{-1/2}
              \sum_{k=0}^K C_k(p_u,1/2)a_u^{-k}\right\},\\
 f_K(u)&=S_{m(u)}(u)+R_{m(u),K}(u),\\
 Q_{X,K}(u)&=S_{m(u)}(u)-S_N(u)+R_{m(u),K}(u).
 \label{eq:negative-initial-data}
\end{align}
The cutoff transfer is exact. Eventually there is at most one adjacent
mode; its initial-data partition is fixed independently of \(z,T\).
No floor is differentiated. The imported bounds at \(\sigma=1/2\) are
\begin{gather}
 |C_0|\le\tfrac12,\qquad
 |C_k|\le\frac{3\sqrt2}{2\pi}\frac{\Gamma(k/2)}{2^k},\quad k\ge1,
 \label{eq:negative-coefficients}\\
 |H(0,u)-f_K(u)|\le C G(u)a_u^{-1/2}
 \Gamma((K+1)/2)(1.1/a_u)^{K+1}.
 \label{eq:negative-high-order-remainder}
\end{gather}
For \(K\le a_{\min}^2\), odd/even ratio summation gives
\(\sum_{k=1}^K|C_k|a_u^{-k}\le C/a_u\), without hiding a
\(K\)-dependent exponential. In particular
\(|Q_{X,K}(u)|\le C G(u)X^{-1/4}\).
Exact \(\Mc\) is retained, rather than a fixed-order Stirling factor.

Already the bounded heat transform of \(f_K\) is a finite representation
at the requested outer precision. To replace its arithmetic integrals
by explicit Gaussian moments, set, for \(\epsilon=\pm1\), the following
center-frozen data. Here \(w_\epsilon,h_{\epsilon n}\) are evaluated
at \(X\), in contrast to the target-dependent functions in
\eqref{eq:first-mode-data}:
\begin{gather}
 w_\epsilon=(1+\epsilon iX)/2,\quad
 h_{\epsilon n}=\Mc(w_\epsilon)n^{-w_\epsilon},\quad
 \lambda_{\epsilon n}=\frac{\epsilon i}{2}
                    [\alpha(w_\epsilon)-\log n],\\
 A_\epsilon(q)=\frac{\Mc(w_\epsilon+\epsilon iq/2)}{\Mc(w_\epsilon)}
       e^{-\epsilon i\alpha(w_\epsilon)q/2},\qquad
 a_{\epsilon k}=A_\epsilon^{(k)}(0)/k!.
 \label{eq:negative-common-curvature}
\end{gather}
The factor \(A_\epsilon\) is independent of \(n\).
Sectorial gamma bounds give \(|\log A_\epsilon(q)|\le C|q|^2/X\)
on \(|q|\le R\). For a sufficiently large onset, Cauchy's formula gives
\[
 |a_{\epsilon k}|\le2R^{-k},\qquad
 \left|A_\epsilon(q)-\sum_{k=0}^Ja_{\epsilon k}q^k\right|
 \le4(V/R)^{J+1},\qquad |q|\le V.
\]
Thus the entire arithmetic heat sum is approximated by the finite moments
\begin{equation}
 \mathcal M_{X,N,J}(T,z)=\sum_{\epsilon=\pm1}\sum_{n\le N}
 h_{\epsilon n}\sum_{k=0}^J a_{\epsilon k}
 \left.\partial_\lambda^k e^{\lambda(z-X)+T\lambda^2}
 \right|_{\lambda=\lambda_{\epsilon n}}.
 \label{eq:negative-moment-sum}
\end{equation}
The final representation is
\begin{equation}
 \widehat F_X(-T,z)=\mathcal M_{X,N,J}(T,z)
 +\frac1{\sqrt{4\pi T}}\int_{I_X}
       e^{-(z-u)^2/(4T)}Q_{X,K}(u)\dd u.
 \label{eq:negative-finite-representation}
\end{equation}
It retains every main arithmetic mode and \(K+1\) RS terms, with a
bounded correction integral of named functions. It is not a short
elementary cosine or a useful-cost theorem.

\begin{proposition}[Bounded fixed-negative-time comparison]
Uniformly for \(1\le T\le2\), on every full fixed buffered chart,
after a finite onset depending on those buffers,
\begin{align}
 |\partial_z^j(H-\widehat F_X)(-T,z)|
       &\le C_j\mathcal E_o(-T,z)e^{-L^2/4},&&0\le j\le3,
       \label{eq:negative-outer-jets}\\
 |\partial_t(H-\widehat F_X)(-T,z)|
       &\le C\mathcal E_o(-T,z)e^{-L^2/4}.
       \label{eq:negative-outer-time}
\end{align}
This is proved in-note from named classical inputs, with same-model
internal assessment only. No numerical onset is supplied.
\end{proposition}

For the proof, use \(\mathcal E_o(-T,z)\ge cG(X)e^{-TL^2/16}\).
The integrated initial remainder has outer-relative logarithm at most
\(-3L^2/8+(L\log L)/2+O(L)\); this includes
\(\Gamma((K+1)/2)\). The complete arithmetic population is bounded
by \(CG(X)X^{1/4}\), with uniformly bounded real exponential tilts.
The common curvature error has logarithm at most
\(-7L^2/8+3L\log L+O(L)\) after outer conversion.
The artificial polynomial tails are bounded by that population times
\(e^{-V^2/64}+e^{-R^2/128}\); the Gaussian dominates the degree
\(J=O(\log X)\) on the whole far ray.

The genuine tails on \(V<|u-X|\le X/2\) retain
\(X^P e^{-\pi X/8+\pi|u-X|/8}\), for fixed \(P\), before Gaussian
integration. Their outer-relative leading exponent is \(-7L^2/8+O(L)\).
Only beyond \(|u-X|=X/2\) is the global source bound used.
Spatial kernel derivatives through order three add fixed polynomial
moments on this interval. Finally both terms of
\eqref{eq:negative-finite-representation} solve \(\partial_TF=F_{zz}\)
exactly, with all parameters and endpoints frozen. The time comparison
therefore follows from the proved second spatial derivative, without
differentiating an error allowance. Note05 supplies the full estimates
and eventual absorption gates.

\subsection{The shrinking-to-fixed bridge and a paid time join}
Ansatz Note06 \cite{ansatz06} extends the unchanged finite model
\eqref{eq:negative-finite-representation} to \(4/L\le T\le2\).
The new issue is the complex kernel as its lower time endpoint shrinks.
For fixed \(B\ge r+b+1\), set
\[
 D_j(T,B)=T^{-j/2}(1+B/\sqrt T)^j,\qquad D_0=1.
\]
Gaussian rescaling, after Hermite differentiation of the kernel, gives
\begin{equation}
 \int_{\R}\left|\partial_z^j
       \frac{e^{-(z-X-q)^2/(4T)}}{\sqrt{4\pi T}}\right|
       e^{c|q|}\dd q
 \le C_j e^{b^2/(4T)}D_j(T,B),\qquad j\le3,\quad0<T\le2.
 \label{eq:bridge-kernel}
\end{equation}
The original \(T^{-1/2}\) normalization cancels only at this rescaling
step. On \(T\ge4/L\), the new transverse factor is at most
\(e^{b^2L/16}\) and \(D_j\le C_j(1+L)^j\).
These losses add \(O_{r,b}(L+\log L)\) to the previously paid
logarithmic error budgets; they do not consume their quadratic slack.

In particular the RS remainder's outer-relative logarithm is at most
\(-3L^2/8+(L\log L)/2+O_{r,b}(L)+j\log(1+L)\).
The full arithmetic curvature term has
\(-7L^2/8+3L\log L+O_{r,b}(L)+j\log(1+L)\).
The genuine intermediate tail retains its height exponential and costs
\[
 C_jX^{P-7/4}D_j(T,B)
 e^{TL^2/16+b^2/(4T)-V^2/(16T)}.
\]
The distant true tail and both artificial polynomial tails are also
paid in Note06, uniformly in the moving time domain. All parameters
and exact cutoff transfers remain frozen. Consequently
\begin{equation}
 |\partial_z^j(H-\widehat F_X)(-T,z)|
 \le C_j\mathcal E_o e^{-L^2/4},\quad j\le3,\qquad
 |\partial_t(H-\widehat F_X)|\le C\mathcal E_o e^{-L^2/4}
 \label{eq:bridge-theorem}
\end{equation}
on \(4/L\le T\le2\). The exact heat identity of this frozen model
is unchanged; its direct time comparison consumes the second spatial
comparison. No constraint defining the moving domain is differentiated.

To join through zero, let \(F_c=F_N\) be exactly Note04's crossover
model, \(F_b=\widehat F_X\), and choose a fixed smooth scalar step
\(\chi(s)\) equal to zero for \(s\le4\) and one for \(s\ge8\), flat
at both endpoints. Use \(F_c\) for \(TL\le4\), \(F_b\) for
\(TL\ge8\), and on the paid overlap set
\begin{equation}
 F_{\rm j}(-T,z)=(1-\chi(TL))F_c(-T,z)+\chi(TL)F_b(-T,z).
 \label{eq:negative-time-join}
\end{equation}
This piecewise definition avoids evaluating \(F_b\) at \(T=0\).
Near zero it inherits Note04's through-zero regularity. Since \(L\)
is frozen at X, the weight has no spatial derivatives.

Write \(a=1-\chi(TL)\), \(\theta=\chi(TL)\),
\(\delta=X^{-3/4}e^{TL^2/16}\) only where \(TL\le8\), and
\(\rho=e^{-L^2/4}\) only where \(TL\ge4\).
Absent branch terms below are omitted. The more precise allowances are
\begin{align}
 |\partial_z^j(H-F_{\rm j})|
  &\le C_j\mathcal E_o(a\delta\Lambda_X^j+\theta\rho),\quad j\le3,
       \label{eq:joined-weighted-jets}\\
 |\partial_t(H-F_{\rm j})|
  &\le C\mathcal E_o\{a\delta\Lambda_X^2+\theta\rho
          +L|\chi'(TL)|(\delta+\rho)\}.
       \label{eq:joined-weighted-time}
\end{align}
The extra time term follows from \(\partial_t\chi(TL)=-L\chi'(TL)\)
and \(|F_b-F_c|\le C\mathcal E_o(\delta+\rho)\), proved by comparing
both models with H. It cannot be dropped.

With \(R_c=(\partial_t+\partial_{zz})F_c\), the actual signed
heat residual is
\begin{align}
 R_{\rm j}&=aR_c-L\chi'(TL)(F_b-F_c),
             \label{eq:joined-residual}\\
 |R_{\rm j}|&\le C\mathcal E_o\{a\delta\Lambda_X
                  +L|\chi'(TL)|(\delta+\rho)\}.
\end{align}
Only the derivative-of-step term is confined to the overlap. The total
residual also contains the crossover residual near zero, and vanishes
for \(TL\ge8\). Thus the complete negative interval has the uniform
chartwise estimates
\begin{align}
 |\partial_z^j(H-F_{\rm j})|&\le C_j\mathcal E_oX^{-1/4}\Lambda_X^j,
                           &&0\le j\le3,\quad -2\le t\le0,
                       \label{eq:negative-interval-jets}\\
 |\partial_t(H-F_{\rm j})|&\le C\mathcal E_oX^{-1/4}\Lambda_X^2,
 & |R_{\rm j}|&\le C\mathcal E_oX^{-1/4}\Lambda_X.
                       \label{eq:negative-interval-time}
\end{align}
For \(TL\ge8\), the stronger exponential field/jet/time bound holds
and the residual is exactly zero. These are proved-in-note estimates
from named inputs, with same-model internal assessment only and no
evaluated onset. Physical spatial joins and nonlinear observable
transfer remain separate obligations.


\subsection{Remaining negative-time limits}
The joined result retains its original \(-2\le t\le0\) scope,
with actual time defect. Proposition~\ref{prop:compact-negative-family}
now pays arbitrary declared fixed compact negative intervals bounded
away from zero at the natural gamma scale, with explicit endpoint
dependence. It does not itself join those models through zero or
extend the older joined outer-scale theorem. Integral-free short
formulas, outer-scale resummation, useful coefficient/precision/quadrature
cost and joint negative-time limits remain open. The RS integral and
growing outer-relative allowance are retained.

The negative endpoint Gaussian
\begin{equation}
F_-(t,z)=\frac{\Phi(0)\sqrt\pi}{2\sqrt T}
                       e^{-z^2/(4T)},\qquad t=-T,
\label{eq:negative-gaussian}
\end{equation}
is an inherited large-\(T\) approximation, with strong absolute density
comparisons in the parent construction. This does not make it a
carrier-relative large-\(X\) approximation at fixed \(T\). An absolute
error can exceed an exponentially small leading field. Growing negative
times and a proved overlap with \eqref{eq:negative-gaussian} remain
separate; their scales are not inferred from a positive-time saddle or
a small heat residual. Physical observable transfer and all-time
assembly also remain open.


\section{Positive-time continuation and joint large-height limits}
\label{sec:positive-coverage}
The signed crossover reaches only a shrinking positive-time window.
The published small-positive-time input on \(0<t\le1/2\) and
the parent comparison on \(t\ge32\) are distinct imports; their
combination does not fill \(1/2<t<32\) or pay uniform compact-time
joins. Closing this range with canonical arithmetic and gamma
estimates is a primary task. The residual hierarchy remains useful,
but positive-time error propagation requires analytic source control
as described in Section~\ref{sec:evolution}.

\subsection{Full saddle and growing positive time}
\label{sec:joint}
When the fixed-time cosine ceases to have a controlled error, keep the
parent full gamma saddle as the reference:
\begin{align}
p_t(w)&=w+\frac t2\alpha(p_t(w)),&
B_t(w)&=1-\frac t2\alpha'(p_t(w)),\\
\Sc_t(w)&=\frac{\Mc(p_t(w))}
 {\sqrt{B_t(w)}}\exp\left\{-\frac{(p_t(w)-w)^2}{t}\right\},\\
P_*(t,z)&=\Sc_t(s)+\Sc_t(1-s). \label{eq:full-saddle}
\end{align}
Use the specified large branch and the square root near one.
The parent controls distinguish static error, evolution error and heat
residual. With
\[
N_*=1+X+t\log(2+t+X),
\]
their characteristic scales are respectively
\(t^{3/2}/N_*^2\), that quantity times \(\log^2(2+t+X)\), and
\(t/N_*^3\). Their actual domains and finite constants stay attached
to the original results.

An optional simplification replaces the digamma saddle by
\begin{equation}
p_0=w+\frac t4\Log\frac{p_0}{2\pi}
 =-\frac t4W_{\mathcal B}\left(-\frac{8\pi}{t}e^{-4w/t}\right).
\label{eq:W-saddle}
\end{equation}
Set
\begin{align}
\Mc_{\rm el}(p)&=\frac{\sqrt\pi}{8}p^{1/2}(p-1)
 \exp\left\{\frac p2\left(\Log\frac p{2\pi}-1\right)\right\},\\
\mathcal C_W(t,w)&=
\frac{\Mc_{\rm el}(p_0)e^{-(p_0-w)^2/t}}
     {\sqrt{1-t/(4p_0)}},&
F_W&=\mathcal C_W(t,s)+\mathcal C_W(t,1-s).
\label{eq:W-field}
\end{align}
This is a candidate. The branch is continued in \(w\) and reflected
conjugately; a fixed Lambert branch index on a winding exponential
argument is not a valid global prescription \cite{dlmf-W}.
The \(W\) replacement needs separate field, phase, derivative and
heat-residual bounds. It is secondary to the arithmetic and fixed-time
outer work.

\section{From field residuals to physical evolution}
\label{sec:observable-evolution}
The physical evolution equations already accommodate an approximate
field. Parent Notes31 and 34 \cite{parent31,parent34}, summarized in
Ansatz Note08, supply the following identities for a spatially
holomorphic \(F\). Retain the physical \(q\) and fixed \(K\), put
\(D_F=|F|^2+q\), and define
\begin{equation}
 \Xi_F=\frac{2\Re(\overline F\,\mathfrak r_F)}{D_F},\qquad
 \mathcal B_F=\frac{q_t+(q_{xx}-q_{yy})/2}{D_F},\qquad
 \square=\partial_x^2-\partial_y^2.
 \label{eq:observable-forcing}
\end{equation}
The actual covariance remains \eqref{eq:covariance}. Direct
differentiation and fixed spatial convolution give
\begin{equation}
 (U_F)_t=-\tfrac12\{\square U_F+(U_F)_x^2-(U_F)_y^2+Q_F\}
           +K*\mathcal B_F+K*\Xi_F.
 \label{eq:forced-potential}
\end{equation}
Define the regulator curvature by the identity
\(\Delta S_F=4\pi(\rho_F+C_F)\), and write
\(\overline C_F=K*C_F\). The corresponding density equation is
\begin{equation}
\begin{split}
 (R_F)_t={}&-\tfrac12\square R_F
 -\frac1{8\pi}\Delta\{(U_F)_x^2-(U_F)_y^2+Q_F\}\\
 &+\frac1{4\pi}\Delta(K*\mathcal B_F)
   -\partial_t\overline C_F-\tfrac12\square\overline C_F
   +\frac1{4\pi}\Delta(K*\Xi_F).
\end{split}
\label{eq:forced-density}
\end{equation}
Thus the ansatz residual is an explicit forcing, rather than an
unaccounted error in a putative closed density equation.

In the current balance its additional forcing is
\begin{equation}
 \mathcal D_{\mathfrak r_F}
 =\frac{2q\Re((\mathfrak r_F)_z\overline{F_z})}{\pi D_F^2}
  -4\rho_F\Re\frac{\mathfrak r_F\overline F}{D_F}.
 \label{eq:forced-current}
\end{equation}
It is added to the existing variable-regulator source in
\((R_F)_t+\operatorname{div}\mathbf J_F\), after convolution
by \(K\), with the parent's complex-to-real current convention.
The equations for \(H\) have \(\mathfrak r_H=0\). Subtracting
them retains differences of the already defined covariance and
regulator terms; it does not define a closure to make them agree.

For the models \eqref{eq:first-negative-model} and
\eqref{eq:compact-negative-model}, insert their respective actual
residuals \eqref{eq:first-negative-residual} and
\eqref{eq:compact-negative-bounds}. In parent notation PE37--PE39,
\(Z_0=\mathfrak r_F\overline F/D_F\); the potential adds
\(2\Re K*Z_0\), density adds
\(\Delta\Re K*Z_0/(2\pi)\), and current balance adds
\(K*\mathcal D_{\mathfrak r_F}\). The numerator bounds are paid,
but the first unpaid physical consumer is a regulator-uniform
weak/collar comparison for this model and \(H\), sufficient for the
parent PE14/PE35 estimates and the actual signed \(Q\) comparison.
No simple-root collar or outer-relative smallness has been acquired
here; the positive-time parent transfer theorem cannot simply be invoked.

For negative time the physical regulator depends on time alone, so
\(C_F=0\), but \(\mathcal B_F=q_t/D_F\) generally remains.
The spatial regulator curvature returns near its switch. A spatial
join that breaks holomorphicity requires the further compatibility
and curvature defects in parent Note42 \cite{parent42}.

These formulas identify the consumers of an asymptotic error bound.
They do not by themselves control those consumers. For example
\( |\Xi_F|\le |\mathfrak r_F|/\sqrt q\) can be poor when the
regulator is small; derivatives and narrow zero collars require
additional estimates. The parent weak-transfer methods are available,
but their hypotheses and constants must be paid for the new ansatze.
The priority is to improve the field, its spatial jets, its direct
time error and its residual together, then prove the physical transfer
on the intended time regions. A coarse density tending to zero does
not resolve the underlying normalized field.

\section{The regulator transition near time 64}
\label{sec:regulator}
The regulator transition remains an important part of the all-time
observable program. Its fine-scale analysis is supporting work after
the immediate field-error improvements. The parent physical protocol
is retained exactly. With
\(C=\Phi(0)\sqrt\pi/2\), \(\eta>0\), and the positive saddle quantities
defined in Appendix~\ref{app:catalogue}, put
\[
\epsilon_b(t)=\eta C(1+t^2)^{-1/4},\qquad
\epsilon_+(t,x)=M_0(t)\kappa(t)^{-1}e^{-x^2/(2\kappa(t))}.
\]
For \(64<t<128\), let \(v=(t-64)/64\) and
\[
\chi(t)=\frac{e^{-1/v}}{e^{-1/v}+e^{-1/(1-v)}}.
\]
Extend \(\chi\) by zero before 64 and one after 128. Then
\begin{equation}
\epsg(t,x)=
\begin{cases}
\epsilon_b(t),&t\le64,\\
\epsilon_b(t)^{1-\chi(t)}\epsilon_+(t,x)^{\chi(t)},&t>64.
\end{cases}                                      \label{eq:regulator}
\end{equation}
These numbers specify the measurement protocol, not intrinsic critical
times of \(H\).

Use the \emph{one-arm} linear amplitude
\(A_X(t)=|\Gc_t((1+iX)/2)|\), its exact local phase derivative
\(\omega(t,X)\), and
\begin{equation}
b(t,X)=\log\frac{\epsg(t,X)}{A_X(t)},\qquad
d(t,X)=\frac{b(t,X)_+}{\omega(t,X)}.
\end{equation}
The inherited coarse model is
\begin{equation}
R_{\rm layer}(t,X,y)=\frac{\omega(t,X)}{2\pi}
              [k_h(y-d)+k_h(y+d)].                 \label{eq:layers}
\end{equation}
On its declared domains its absolute error is \(O(1)\), while
\(\omega\sim\frac14\log(X/(4\pi))\). These are layers of regularized
density, not nonreal zeros.

The exact model center is \(b(t_c(X),X)=0\).
The parent gives
\begin{equation}
t_c-64=\frac{64}{\log X+c_0+O(1/\log X)},\qquad
c_0=\log\frac{4e}{\pi\kappa(64)}.
\label{eq:center}
\end{equation}
Its rounding scale is much smaller:
\begin{equation}
t_X(\sigma)=t_c(X)-\frac{128\sigma}{\pi X\log X},\qquad
\frac{R_H(t_X(\sigma),X,y)}{\log X}
=\frac{k_h(y-\sigma_+)+k_h(y+\sigma_+)}{8\pi}
 +O(1/\log X).                                    \label{eq:rounding}
\end{equation}
The symbol \(\sigma\) here is unrelated to the arithmetic coordinate
\(\tau=t\ell_X\).

Replacing \(t_c\) by a short inverse-logarithmic expansion is not justified
at the scale \eqref{eq:rounding}. For the next-order transition, even
the factor two between \(A_{\mathrm o}\) and \(A_X\) matters:
it changes \(b\) by a nonzero constant. Substituting an elementary or
full-saddle amplitude for \(A_X\) needs an explicit recentering estimate.
Direct time derivatives of this transition must be derived separately,
not obtained by differentiating an \(O(1)\) density allowance.

At fixed time the parent protocol has the nonuniform limit
\[
\frac{R_H(t,X,y)}{\log X}\longrightarrow
\begin{cases}0,&t\le64,\\ k_h(y)/(4\pi),&t>64,\end{cases}
\qquad |y|<h.
\]
Before the switch, a vanishing coarse density does not resolve the
normalized field or its arithmetic oscillations. Both questions remain
in the high-height program.

\clearpage
\section{All-time coverage and quantitative acceptance}
\label{sec:map}
The following map records the domains actually supported by the
current sources and the remaining field-estimate tasks. It is not a
completed covering theorem. Here \(L=\log(X/(4\pi))\).
\begingroup\small
\begin{longtable}{@{}>{\raggedright\arraybackslash}p{.20\textwidth}>{\raggedright\arraybackslash}p{.36\textwidth}>{\raggedright\arraybackslash}p{.36\textwidth}@{}}
\toprule Region & Present evidence & Primary remaining estimate\\\midrule
\endhead
\(t=0\) and \(|t|L\le8\) &
Note04: corrected fixed-cutoff complex comparison, spatial jets
through three and direct first time error on full fixed buffers. &
Higher controlled orders when needed; canonical initial matching
for evolution corrections and compatibility on overlaps.\\
\(-2\le t\le0\) &
Note06: chartwise joined field and derivative estimates, with the
actual time-cutoff residual retained. &
Pay a declared join to the separate compact-away-zero family,
retaining the component and cutoff residuals; no enlarged join is
claimed here.\\
\(1\le T=-t\le2\) &
Note10: fixed-order retained-integral formula at \(|G|X^{-5/4}\)
field/jet/time error and \(|G|X^{-7/4}\) residual. Note05's stronger
growing-order outer model remains. &
Integral-free simplification and outer-scale resummation: allowance
\(X^{-5/4}e^{TL^2/16}\) still grows. Physical weak/collar transfer unpaid.\\
Declared fixed \(0<\tau\le T\le B\) &
Note12: exact-gamma retained-integral model with explicit endpoint
factors/onsets, \(|G|X^{-5/4}\) field/jets/time and
\(|G|X^{-7/4}\) residual. &
No zero/growing-time substitution. Constants deteriorate at endpoints;
outer allowance still grows. Physical transfer remains unpaid.\\
Fixed positive time &
Distinct published \(0<t\le1/2\) and parent \(t\ge32\) inputs;
Note03's corrected buffered theorem on \([32,33]\). &
Pay the intermediate positive range \(1/2<t<32\) and uniform
compact-time overlaps; retain arithmetic terms whenever larger than
the requested correction.\\
Joint \(t,X\to+\infty\) &
Inherited full gamma saddle with its stated domains and constants. &
Uniform large-height reduction and matching; a Lambert \(W\)
replacement only if it gives a proved gain.\\
Joint \(T=-t,X\to\infty\) &
Inherited large-\(T\) Gaussian endpoint with absolute comparisons. &
Identify scaling regions and prove high-height overlaps; do not
substitute it for fixed-negative-time relative asymptotics.\\
Regulator transition near time 64 &
Inherited two-layer density profile with \(O(1)\) absolute error
and an exact implicit center. &
After the primary field work: one-arm amplitude comparison,
recentering at the rounding scale, and direct dynamic transfer.\\
\bottomrule
\end{longtable}
\endgroup

For each proposed model record a positive carrier envelope
\(\mathcal A(t,X)\), independent of the zeros of its signed sum.
An accepted estimate must state the time domain, height onset, full
complex buffer and parameter dependence, and supply separately:
\begin{enumerate}
\item canonical initial mismatch, including the gamma factor,
all retained arithmetic modes and the cutoff correction;
\item field and spatial errors through the orders needed by the
consumer, currently at least \(j\le3\);
\item the direct first time error and actual residual, with every
time or spatial joining defect;
\item the propagation estimate or direct canonical comparison that
turns these data into a solution-error bound;
\item overlap errors and, when claimed, transfer to the physical
\(U,W,R,J,Q\), regulator terms and remaining parent tuple.
\end{enumerate}
Keep natural-gamma-scale and outer-carrier-scale allowances distinct.
An improvement in an absolute norm can still leave the desired
relative estimate unproved. A residual identity alone does not
establish a carrier-relative field error. Likewise, an algebraically
shorter formula is not an error improvement, and a longer exact heat
representation need not have useful evaluation cost.

The first bounded negative-time experiment now attains
\(C|G|X^{-5/4}=o(|G|X^{-3/4})\) with the required jets and time
error, in the fixed-order retained-RS-integral class. The initial,
population, kernel and tail estimates are separate from residual
cancellation. Its outer allowance \eqref{eq:first-remaining-outer-loss}
still fails; no actual error lower bound is inferred. The fixed
compact-negative extension is now paid by
Proposition~\ref{prop:compact-negative-family}, with the full explicit
ledger in Appendix~\ref{app:compact-endpoints}. Through-zero assembly,
the missing positive range and matching joint limits remain primary tasks.

\section{Scope, controls, and continuation priorities}
These are approximation models for one canonical source. They are
not an RH criterion, a simplicity theorem or a collision exclusion.
Heat evolution, Duhamel's identity and the modulation equations also
hold for the project's collision controls and exclude none.
The source distinctions are analytical:
\begin{itemize}
\item Euler-factor perturbations change the ordinary Dirichlet
coefficients. Finite-prime agreement does not recover the full
canonical arithmetic expansion.
\item Finite weighted-square and theta-cutoff controls have different
complete coefficients and tails; selected endpoint equalities do not
supply the canonical completion and RS identity.
\item Reflected super-Gaussian, smooth-shape, positive-primitive,
counter-kernel and fixed-event controls lack the same complete
gamma/arithmetic identity. Their shared positivity or heat equation
cannot pay a missing canonical remainder.
\end{itemize}
These premise distinctions do not themselves prove a geometric
conclusion. The source-specific quantitative comparison must do the
work. Consult the repository result and avoidance indexes
\cite{controls,avoidance}, especially the cautions on shared premises,
separate bounds that lose cancellation, same-model review and statuses
stronger than their evidence.

The continuation order is now:
\begin{enumerate}
\item \textbf{Controlled large-height field asymptotics across time.}
The first explicit evolution correction at intermediate negative
time is integrated in Section~\ref{sec:negative-first-correction},
with its two accuracy scales and retained integral. The declared
fixed compact-negative family is now integrated in
Section~\ref{sec:compact-negative-family}. Next fill the positive-time
coverage gap and pay the relevant compact-family and through-zero
overlaps. Derive joint positive and negative regimes with their own
uniform bounds; their scope is part of this first priority.
\item \textbf{Transfer to the physical evolution and assemble regions.}
Track the residual forcing, actual covariance, fixed regulator and
joining errors. Transfer may be developed on completed domains without
claiming a global result prematurely.
\item \textbf{Resolve the fine regulator transition near time 64.}
Retain the prepared one-arm/center-comparison task, the exact implicit
center and the much smaller rounding scale. It is a supporting
workstream, not the next task ahead of the asymptotic error estimates.
\item \textbf{Simplify secondary representations and assess cost.}
Lambert \(W\), the central Gaussian-polynomial model and evaluation
efficiency are useful when they preserve or strengthen paid estimates.
\end{enumerate}
The accompanying overview and progress register implement this order.
HDA-A02-AN03 and AN05 close their bounded analytical results; AN04
and AN06 record the respective assessed-result integrations. The
broader HDA-A02 remains Partial. The next prepared bounded field task
is HDA-A08-AN01, Analysis---Sol, a canonical buffered estimate on the
first missing positive-time slab \([1/2,1]\). It is prepared only;
no positive-time extension follows from the negative-time integration. The regulator packet HDA-A03-AN01 remains retained for later
work. Parent HDA-17 Partial, HDA-18 Gated and HDA-20 unfulfilled are
unchanged. No independent review or physical-transfer completion is
inferred from integration or compilation.

\clearpage
\appendix
\section{Endpoint factors and onsets for the compact-negative family}
\label{app:compact-endpoints}
This ledger supplies the explicit dependence in
Proposition~\ref{prop:compact-negative-family}; it is not a new
zero-time or growing-time theorem. Follow Note12 CF4--CF30
\cite{ansatz12}, with target \(r=\ell+2\), \(b=1+h+2\), enlarged
\(R=r+1\), \(Y=b+1\), and radius-\(1/2\) Cauchy disks.
The \(R\) here is a fixed chart radius, distinct from the
\(X^{1/3}\) Taylor radius in \eqref{eq:negative-orders}.
Set \(a_0=\pi/8\). Fixed classical constants
\(C_0,D,g_1,g_2,g_3,g_->0\) and \(X_\Gamma(R,Y)\), independent
of \(\tau,B\), are chosen to satisfy the gamma/RS bounds
\begin{gather}
 G(X)\ge g_-X^{7/4}e^{-\pi X/8},\quad |G(z)|\asymp_{R,Y}G(X),
 \quad G(u)/G(X)\le C_0e^{|u-X|},\\
 \beta=-a_0+\epsilon ir_n/2+\delta,\quad |\delta|\le D/X,
 \quad |\gamma|\le g_1/X,\quad |\beta''|\le g_2/X^2,
 \quad |\beta'''|\le g_3/X^3.
 \label{eq:compact-classical-data}
\end{gather}
The initial \(K=1\) mismatch is bounded by \(C_0G(u)X^{-5/4}\) and
\(|Q_{X,1}(u)|\le C_0G(u)X^{-1/4}\) on the frozen window.
The classical constants and onset are analytical, not numerically
instantiated. Real initial boundedness is
\(M_H=\int_0^\infty|\Phi(v)|\dd v\).

\subsection{Kernel and complete-population factors}
For integers \(m,d\ge0\), define
\begin{align}
 \mathcal M_m&=\sqrt2e^{R+Y^2/(4\tau)+2B}
 \sum_{k=0}^m\binom mk Y^{m-k}(8B)^{k/2}
                  \frac{\Gamma((k+1)/2)}{\sqrt\pi},\\
 \mathcal T_m&=\frac{\mathcal M_m}{2\tau}
                     +\frac{\mathcal M_{m+2}}{4\tau^2}.
 \label{eq:compact-kernel-factors}
\end{align}
They majorize, respectively,
\(\int|k_T(z-u)|e^{|u-X|}|u-z|^m\dd u\) and the same integral
with \(|\partial_Tk_T|\). Indeed \(|u-X|\le|u-x|+R\),
\(|u-z|\le|u-x|+Y\), and
\(|u-x|\le (u-x)^2/(8T)+2T\). The resulting real Gaussian
moments are exactly the Gamma factors displayed above. Time inserts
\((z-u)^2/(4T^2)-1/(2T)\). Thus the transverse exponential and
inverse-time costs are explicit.

Put \(c_Y=(Y-1)/2\) and
\begin{align}
 W_d(\tau)&=2^d+\frac{3^d}{2}\sum_{k=0}^d\binom dk
          (8/\tau)^{(k+1)/2}\Gamma((k+1)/2),\\
 \Pi_d&=C_0^{d+1}e^{Ba_0^2+1+(Y+1)/2}
                 e^{2(c_Y)_+^2/\tau}W_d(\tau).
 \label{eq:compact-population-factors}
\end{align}
Under the slope gate below,
\begin{equation}
 \sum_{\epsilon,n\le N}|h_{\epsilon n}e^{T\beta_{\epsilon n}^2}|
       (1+|\beta_{\epsilon n}|)^d
       \le \Pi_dG(X)X^{1/4}.
 \label{eq:compact-population}
\end{equation}
To see this, \(|e^{T\beta^2}|\le e^{Ba_0^2+1}e^{-\tau r_n^2/4}\)
and \(|h(z)|\le C_0G(X)n^{-1/2}e^{Yr_n/2}\).
Each shell \(j\le r_n<j+1\) has at most \(2a_Xe^{-j}\)
integers. Use
\(c_Yj-\tau j^2/4\le2(c_Y)_+^2/\tau-\tau j^2/8\).
For \(j\ge1\), \((j+2)^d\le3^dj^d\), and comparison with the
integral of \((x+1)^de^{-\tau x^2/8}\) over \([j-1,j]\)
gives \(W_d\). All modes are paid before estimating the signed sum.
The first correction's absolute population is bounded by
\(g_1(B\Pi_d+2B^2\Pi_{d+2})G(X)X^{-3/4}\) after the same moment
weight. This is not a leading asymptotic for \(H\).

The exact residual expansion of \eqref{eq:first-exact-residual} gives
\(|(\mathcal A_T-\mathcal A_{zz})/(he^{T\beta^2})|
\le X^{-2}(P_2+X^{-1}P_3)(1+|\beta|)^4\), where
\begin{align}
 P_2(B)&=g_1^2(3B+18B^2+8B^3)+4g_2(B+B^2),\\
 P_3(B)&=g_3(B+2B^2)+g_1g_2(18B^2+12B^3)\notag\\
       &\hspace{1em}+g_1^3(6B^2+24B^3+8B^4).
 \label{eq:compact-residual-polynomials}
\end{align}
The factors used in the theorem are
\begin{equation}
 \begin{gathered}
 \mathcal E=3+\mathcal M_0+\mathcal M_3+\mathcal M_4,\qquad
 \mathcal E_t=3+\mathcal T_0+\mathcal T_3+\mathcal T_4,\\
 \mathcal R=\Pi_4\{P_2(B)+P_3(B)\}.
 \end{gathered}
 \label{eq:compact-error-factors}
\end{equation}
The fixed Cauchy/gamma comparison constants are absorbed into
\(C_{r,b,j}\) and \(C_{r,b}\), independently of \(\tau,B\).

\subsection{True and artificial tails}
Put \(v=u-X\); physical regulator \(q\) is unchanged.
Define
\begin{align}
 \mathcal D_m&=\frac{e^{Y^2/(4\tau)}}{\sqrt{4\pi\tau}}
       2^m(64B)^{(m+1)/2}\Gamma((m+1)/2),\\
 \mathcal J_m&=\frac{e^{Y^2/(4\tau)}}{\sqrt{4\pi\tau}}
 \sum_{k=0}^m\binom mk Y^{m-k}(8B)^{(k+1)/2}\Gamma((k+1)/2).
 \label{eq:compact-tail-factors}
\end{align}
For either family let \(A_m^t=A_m/(2\tau)+A_{m+2}/(4\tau^2)\).
If \(V_B\ge\max(2R,R+Y+1,32B)\), the kernel with real tilt obeys
\[
 -\frac{(v-(x-X))^2}{4B}+|v|\le-\frac{v^2}{32B},\qquad
 |u-z|^m\le(2|v|)^m\quad(|v|\ge V_B).
\]
Its tail moments are at most
\(\mathcal D_m e^{-V_B^2/(64B)}\le\mathcal D_m e^{-L^2}\),
with time factor \(\mathcal D_m^t\).
The genuine moderate field bound retains
\(C_0X^Pe^{-\pi X/8+\pi|v|/8}\), with \(P=4\) admissible.
Thus moderate and artificial tail errors, divided by \(G(X)\), are
respectively at most
\begin{gather}
 C_0X^{P-7/4}\mathcal D_0e^{-L^2},\qquad
 C_0X^{1/4}(\mathcal D_0+g_1X^{-1}\mathcal D_2)e^{-L^2}.
 \label{eq:compact-near-tails}
\end{gather}
For time use \(\mathcal D_0^t,\mathcal D_2^t\).
Only for \(|v|>X/2\), with \(X\ge4R\), apply the global bound
\(M_H\). The remote field tail divided by \(G(X)\) is at most
\begin{equation}
 \frac{M_H}{g_-}X^{-7/4}\mathcal J_0
                  e^{\pi X/8-X^2/(128B)},
 \label{eq:compact-remote-tail}
\end{equation}
with \(\mathcal J_0^t\) for time. No exact gamma or analytic RS
function is continued through remote poles for a polynomial completion.

Set
\begin{align}
 A_{\rm mod}&=C_0\max(\mathcal D_0,\mathcal D_0^t),\\
 A_{\rm art}&=C_0\max(\mathcal D_0+g_1\mathcal D_2,
                       \mathcal D_0^t+g_1\mathcal D_2^t),\\
 A_{\rm far}&=(M_H/g_-)\max(\mathcal J_0,\mathcal J_0^t).
 \label{eq:compact-tail-absorption}
\end{align}
The following sufficient gates place every field and direct-time tail
below \(G(X)X^{-7/4}\):
\begin{gather}
 L\ge P+\sqrt{\log^+[A_{\rm mod}(4\pi)^P]},\quad
 L\ge2+\sqrt{\log^+[A_{\rm art}(4\pi)^2]},\\
 X\ge32\pi B,\quad X\ge\sqrt{256B\log^+A_{\rm far}},
 \qquad\log^+v=\max(0,\log v).
 \label{eq:compact-tail-onsets}
\end{gather}
They imply \(A_{\rm mod}X^Pe^{-L^2}\le1\),
\(A_{\rm art}X^2e^{-L^2}\le1\), and
\(A_{\rm far}e^{\pi X/8-X^2/(128B)}\le1\).

\subsection{Complete sufficient onset and endpoint limits}
In addition to \(X\ge X_\Gamma(R,Y)\), the real RS onset and
\eqref{eq:compact-tail-onsets}, impose
\begin{gather}
 L\ge16,\quad X\ge\max(1,4R),\quad X-V_B\ge X/2,\quad
 a(X+V_B)-a(X-V_B)<1,\\
 V_B\ge\max(2R,R+Y+1,32B),\qquad
 C_0(V_B+R+Y)^2/X\le1,\qquad D/X\le1-a_0,\\
 B\{(D/X)(2a_0+L/2)+D^2/X^2\}\le1,\qquad
 C_0B(1+L)/X\le1,\\
 2g_1\{B+2B^2[a_0+1+(\log2)/2]^2\}\le X.
 \label{eq:compact-local-onsets}
\end{gather}
Here \(a(u)=\sqrt{u/(4\pi)}\). The last gate ensures
\(|1+c|\ge1/2\) in the near-cutoff band for its lower absolute
population comparison; the upper error theorem does not need that
lower comparison. Every gate is eventual for fixed \(\tau,B,r,b\).
They specify a sufficient analytical onset, not a certified numerical
height or an endpoint-independent \(X_0\).

As \(\tau\downarrow0\), the kernel/tail factors contain
\(e^{Y^2/(4\tau)}\), with \(\tau^{-1},\tau^{-2}\) in direct
time estimates. The population has
\(e^{(Y-1)_+^2/(2\tau)}\) and powers through
\(\tau^{-(d+1)/2}\). These displayed majorants need not be optimal.
For increasing fixed \(B\), the kernel factor has \(e^{2B}\)
times finite moments, population has \(e^{\pi^2B/64}\), and
\(P_2,P_3\) have degree at most four. Window and onset costs also
depend on \(B\). None licenses \(\tau=\tau(X)\), \(B=B(X)\),
a uniform zero-time join, or a large-negative-time overlap.

\section{The remaining ansatz catalogue}
\label{app:catalogue}
These formulas preserve the earlier scope while the main effort moves
to large height.

\subsection{A finite-time Gaussian-polynomial core}
Choose \(a>b\) and a real polynomial \(p_m\), fixed once per accuracy.
The nonnegative source \(e^{-au^2}p_m(u^2)^2\) has exact heat field
\begin{equation}
F_{\rm c}(t,z)=[p_m(-\partial_z^2)]^2
 \left[\frac{\sqrt\pi}{2\sqrt{a-t}}e^{-z^2/[4(a-t)]}\right],
 \qquad t<a.                                      \label{eq:core}
\end{equation}
It is a Gaussian times a finite even polynomial. With
\(p_m(v)^2=\sum_{k=0}^{2m}c_kv^k\) and \(a-t=\nu\),
\[
F_{\rm c}=\frac{\sqrt\pi}{2\sqrt\nu}e^{-z^2/(4\nu)}
 \sum_{k=0}^{2m}\frac{(-1)^kc_k}{(4\nu)^k}
 H_{2k}\!\left(\frac z{2\sqrt\nu}\right).
\]
The Hermite convention is the physicists' one.
Gaussian weighted \(L^2\) approximation of
\(e^{au^2/2}\sqrt{\Phi(u)}\), followed by Cauchy--Schwarz, gives
an existence route for absolute field jets on \(t\le b\) and fixed
strips \cite{ansatz01,dlmf-Hermite}. It does not give a useful degree
bound or a high-height relative comparison. This model is not a fixed
source for the positive infinite-time endpoint.

\subsection{The positive inner cosine}
Let \(L\) be the large real solution and define
\begin{align}
4\pi e^{4L}&=2tL+9,&\kappa&=8tL-2t+36,\\
M_0&=2\pi^2e^{tL^2+9L-\pi e^{4L}}\sqrt{2\pi/\kappa}.
\end{align}
The simplified candidate is
\[
F_{\rm in}=M_0e^{-z^2/(2\kappa)}\cos(Lz).
\]
The parent moment carrier is \(M(1-vz^2/2)\cos(\mu z)\).
Replacing its envelope, mass and frequency must be paid separately,
initially on its buffered \(O(\kappa^{1/3})\) inner window.
The exact real saddle can be written
\[
L=-\tfrac14W_{-1}(-(8\pi/t)e^{-18/t})-\frac9{2t}.
\]
A low-order log-log expansion of \(L\) can create an unacceptable
phase error when multiplied by a growing height.

\subsection{Local zero peaks and the clock}
For the local monomial \(F=c(z-a)^m\), freezing the regulator and setting
\(\delta=(\epsg/|c|)^{1/m}\) gives
\[
\rho_{\rm local}=
\frac{m^2\delta^{2m}|z-a|^{2m-2}}
     {\pi(|z-a|^{2m}+\delta^{2m})^2}.
\]
This is exact for the model, not a simplicity assertion for \(H\).

The finite Bernoulli clock formula is exactly the locally finite sum
\begin{align}
d_j&=(j+\tfrac12)\pi/L,\\
R_L&=k_h(y)\sum_j k_\ell(x-d_j),&
J_L&=-\frac2\kappa k_h(y)\sum_jd_jk_\ell(x-d_j).
\end{align}
This pair satisfies continuity. Its continuum leading densities are
\[
R_{\rm in}\simeq Lk_h(y)/\pi,\qquad
R_{\rm out}\simeq \ell_X k_h(y)/(4\pi)
\]
on the respective resolved regimes. Discarding phase is not harmless
for differentiated estimates.

\subsection{Cosine-cell and narrow-filter covariance}
For \(F=M\cos(kz)\), fixed \(k\), constant regulator \(\epsilon\), and
\(M'=k^2M\), let \(\eta_c=\epsilon/M\),
\(a_c=\sinh^2(ky)+\eta_c^2\), and \(d_c=\sqrt{a_c(a_c+1)}\).
The exact phase-cell formulas are
\begin{align}
U_c&=2\log M+2\arsinh\sqrt{a_c}-2\log2,\\
R_c&=\frac{k^2\eta_c^2}{2\pi}
\frac{2a_c(a_c+1)+1-\eta_c^2(2a_c+1)}{d_c^3},\\
Q_c&=4\pi R_c+(U_c)_y^2-4k^2+\frac{4k^2\eta_c^2}{d_c}.
\end{align}
Cell averaging is not the parent compact filter; additional transverse
averaging changes the covariance.
Separately, under the parent's small-width, fixed-regulator and smoothness
hypotheses,
\[
Q\simeq v_t(U_{xt}^2-U_{yt}^2)
       +v_x(U_{xx}^2-U_{xy}^2)+v_y(U_{xy}^2-U_{yy}^2).
\]
This is not an unrestricted closure.

\subsection{Assembly and bounds that need no new ansatz}
A short notation for a joined field is
\[
F=\chi_-F_-+\chi_{\rm c}F_{\rm c}+\chi_+F_+,\qquad
\chi_-+\chi_{\rm c}+\chi_+=1.
\]
Every support, overlap and derivative of these weights must be specified
and paid. Spatial joins can be nonholomorphic. The actual covariance of
the joined field is not the weighted sum of individual covariances.

Finite xi-derivative Taylor formulas remain representations; a formal
heat exponential does not simplify their content. Source derivative
bounds such as parent Lemma B2, tail majorants, moment bounds,
conditioning estimates and evaluation error budgets belong in the
remainder machinery. They need not be simplified to obtain a useful
short solution ansatz.

\begin{thebibliography}{99}
\bibitem{parent} Edward Baker research project,
\emph{Heat-density asymptotics}, maintained parent manuscript,
25 October 2026, through Notes54--55.
Local source: \texttt{../../manuscript/main.tex}.
\bibitem{summary62} Edward Baker research project,
\emph{Chat0021: research results, dependencies and remaining work},
Note62, 25 October 2026; GPT-6.1 Sol, effort ultra.
Local source: parent \texttt{notes/62\_CHAT0021\_RESEARCH\_SUMMARY\_20261025.md}.
\bibitem{ansatz01} Edward Baker research project,
\emph{Simple solution ansatze, separated by regime},
Ansatz Note01, 25 October 2026; GPT-6 Astra, effort max.
\bibitem{ansatz03} Edward Baker research project,
\emph{Outer correction and a bounded signed-time real-axis bridge},
Ansatz Note03, 25 October 2026; GPT-6.1 Sol, effort ultra.
Same-model internal assessment only. Its open complex-buffer gap
was subsequently closed at the stated bounded scope by Note04.
\bibitem{ansatz04} Edward Baker research project,
\emph{Fixed-cutoff complex remainder and buffered signed-time crossover},
Ansatz Note04, 25 October 2026; GPT-6.1 Sol, effort ultra.
Proved in-note from named inputs, same-model internal assessment only.
\bibitem{ansatz05} Edward Baker research project,
\emph{Fixed negative times: a short natural-scale model and an
outer-scale representation}, Ansatz Note05, 25 October 2026;
GPT-6.1 Sol, effort ultra. Proved in-note from named inputs,
same-model internal assessment only; interval \(1\le T\le2\).
\bibitem{ansatz06} Edward Baker research project,
\emph{Negative-time bridge and a paid time join through zero},
Ansatz Note06, 25 October 2026; GPT-6.1 Sol, effort ultra.
Proved in-note from named inputs, same-model internal assessment only;
joined interval \([-2,0]\). Note07 records the chat0001 research summary.
\bibitem{ansatz08} Edward Baker research project,
\emph{Evolution equations, ansatz residuals, and negative-time arithmetic modes},
Ansatz Note08, 25 October 2026; GPT-6 Astra, effort max.
Source-based synthesis and exact identities; no new asymptotic theorem.
Note09 and Assessment06 record this manuscript integration.
\bibitem{ansatz10} Edward Baker research project,
\emph{First negative-time evolution correction with a retained RS integral},
Ansatz Note10, 25 October 2026; GPT-6.1 Sol, effort ultra, chat0003.
Proved in-note from named classical inputs; same-model internal
analytical assessment only. Equations EC4--EC30.
\bibitem{review07} Edward Baker research project,
\emph{Assessment: first negative-time correction and remaining outer loss},
Ansatz Review07, 25 October 2026; GPT-6.1 Sol, effort ultra.
Same-model internal analytical assessment; no specialist review.
\bibitem{ansatz12} Edward Baker research project,
\emph{A fixed compact-negative-time family with endpoint-dependent
allowances}, Ansatz Note12, 25 October 2026; GPT-6.1 Sol, effort ultra.
Proved in-note from the named inputs of Notes05/10; same-model internal
analytical assessment only. Equations CF1--CF32.
\bibitem{review09} Edward Baker research project,
\emph{Assessment: compact-negative-time family and endpoint dependence},
Ansatz Review09, 25 October 2026; GPT-6.1 Sol, effort ultra.
Same-model internal analytical assessment; no specialist review.
\bibitem{parent31} Edward Baker research project,
\emph{Positive all-height evolution and covariance}, parent Note31,
24 October 2026, equations PE33--PE39. Internal research source.
\bibitem{parent34} Edward Baker research project,
\emph{Compact weak transfer and time join}, parent Note34,
24 October 2026, equations CJ25--CJ26. Internal research source.
\bibitem{parent42} Edward Baker research project,
\emph{Gap-localized spatial defects}, parent Note42,
24 October 2026, equations GS7--GS12. Internal research source.
\bibitem{duhamel} Stanford University, Math 220B,
\emph{Heat Equation}, Section 2.7, printed pp.~32--35.
Course handout; author and publication date not assigned here.
\url{https://web.stanford.edu/class/math220b/handouts/heateqn.pdf}.
Local reference catalog records the inspected copy, retrieved
25 October 2026. Only the standard Duhamel identity is imported.
\bibitem{polymath} D. H. J. Polymath,
\emph{Effective approximation of heat flow evolution of the Riemann
xi function, and a new upper bound for the de Bruijn--Newman constant},
arXiv:1904.12438v2 (4 August 2019).
\url{https://arxiv.org/html/1904.12438v2}.
\bibitem{dlmf-gamma} NIST Digital Library of Mathematical Functions,
Section 5.11, equations 5.11.1--5.11.3 and 5.11.9, and sectorial
remainder bounds. \url{https://dlmf.nist.gov/5.11}.
\bibitem{dlmf-rs} NIST DLMF, Section 25.10, equations
25.10.1--25.10.4. \url{https://dlmf.nist.gov/25.10}.
\bibitem{dlmf-W} NIST DLMF, Section 4.13, Lambert W and its branches.
\url{https://dlmf.nist.gov/4.13}.
\bibitem{dlmf-Hermite} NIST DLMF, Section 18.18(i), Hermite expansions
and weighted \(L^2\) convergence. \url{https://dlmf.nist.gov/18.18}.
\bibitem{controls} Repository \emph{Result summaries}, especially
clusters 11--13: corrected Riemann--Siegel interface, collision controls,
and local heat-flow geometry. See the accompanying overview for links.
\bibitem{avoidance} Repository \emph{Patterns to avoid}, especially
01, 02, 05--08, 12, 16 and 19. See the accompanying overview for links.
\end{thebibliography}
\end{document}
